% Section A: manuscript pages 64--90. The overlap with DER_VI-a is cited rather than repeated.
\begin{dcomment}{exposition}{p70:J'ai envie de grouper ce qui a été prouvé}{Three theories within the comparison diagram}
Grothendieck now wants the many presentations to express one intelligible distinction. The proof that supplies it has already been explained in the commentary at manuscript pp.~64, 67, and 69: the path replacement defines a section of the coarse localization, and the two iterated units agree after fine localization. Here the new work is to organize the consequences and identify the image of that section.

For a map $v:X\to Y$ over $S$, the three tests are
\[
\begin{array}{c|l}
 W_S&v\in W\text{ on total categories},\\
 W_S^{\mathrm d}&s\backslash X\longrightarrow s\backslash Y\in W
                    \text{ for every }s,\\
 W_S^{\mathrm g}&X/s\longrightarrow Y/s\in W
                    \text{ for every }s.
\end{array}
\]
The relative Theorem A supplies the inclusions of both fine classes in $W_S$. Thus the identity on $\Cat/S$ gives $\rho$ and $\sigma$; these are quotient functors between homotopy categories, before any assertion about adjoints. In the other direction one must first replace $X$: the functor $C$ sends a coarse equivalence to both kinds of fine equivalence. This is why the inclusion of locally constant objects is represented by a nontrivial construction on total categories.

The replacement by a split fibration does something different. Its fibre at $s$ is $s\backslash X$, so it preserves the right-fine information already present in $X$. Dually, the split opfibration has fibre $X/s$. Neither construction forces its transition maps to become equivalences. The path replacement does force this, and consequently belongs to the upper part of the diagram. Confusing these two purposes would make the distinction between $\HOT_W(S)$ and $\HOT_W^{\mathrm{lc}}(S)$ disappear before it has been proved.

Part~I is therefore a statement about coarse localization of all the presentations. Part~II identifies the appropriate fine localization along each side. Part~III says when a strict fibre is a valid test and when the two fine tests agree. The hypothesis L5 is used at precisely the passage from a coarse equivalence between homotopy fibrations to a fibrewise equivalence. It should not be silently dropped when the diagram is reused. The subsequent remarks at manuscript pp.~81--82 explicitly reopen this dependency.
\end{dcomment}

\begin{dcomment}{exposition}{p71:Sur les six catégories internes du diagramme}{Opposite categories and opposite derivators}
The dual statement is meant to recover the covariant theory from the contravariant one. At the level of the ordinary localized categories used here, its variance can be checked directly. Taking opposites gives a functor
\[
 O_S:\Cat/S\longrightarrow\Cat/S^{\mathrm{op}},
 \qquad (X\to S)\longmapsto(X^{\mathrm{op}}\to S^{\mathrm{op}}),
\]
which sends $v:X\to Y$ to $v^{\mathrm{op}}:X^{\mathrm{op}}\to Y^{\mathrm{op}}$, with the same direction. Moreover
\[
 s\backslash(X^{\mathrm{op}})\cong(X/s)^{\mathrm{op}}.
\]
For the opposite-invariant localizer under consideration, $O_S$ interchanges the left-fine and right-fine classes. Its square is the identity. Consequently the localization definitions of this Scholie give
\[
 (\Cat/S)(W_S^{\mathrm g})^{-1}
       \simeq(\Cat/S^{\mathrm{op}})(W_{S^{\mathrm{op}}}^{\mathrm d})^{-1},
\]
with no opposite on the resulting category.

The editorial insertion on this page asks for an additional outer opposite. That insertion belongs to the different convention of an \emph{opposite derivator}, whose value is $\mathbb D(S^{\mathrm{op}})^{\mathrm{op}}$; it is not compatible with identifying both sides here with the displayed ordinary localizations. We preserve the edition's insertion, and distinguish the conventions rather than altering its words. See also manuscript p.~50, where the same issue first appears. The two diagram functors at manuscript p.~73 confirm the present interpretation: their values lie in $\Hot_W$, with the index category reversed, rather than in $\Hot_W^{\mathrm{op}}$.

A small check prevents a significant error later. At $S=e$ and for ordinary weak homotopy equivalences, both fine localizations are $\Ho(\mathrm{Spaces})$. This category is not equivalent to its opposite: its initial object is strict, since a map to the empty space can only have empty source, whereas its terminal object is not costrict: a map $*\to S^0$ need not be invertible. These properties are preserved by equivalences and exchanged by passage to the opposite. In particular dualizing the present comparison does not reverse the order of the two adjoints: both $\rho\dashv i$ and $\sigma\dashv j$ have the order printed in Part~IV. The later hesitation at manuscript p.~84 comes from precisely this variance distinction.
\end{dcomment}

\begin{dcomment}{development}{p71:Sur les six catégories internes du diagramme}{The groupoid test for collapse of the comparison}
Grothendieck asks in footnote~46 whether the distinction between coarse and fine localization persists for every base not equivalent to a discrete category. A group with one object already shows that the proposed dividing line is too narrow. The relevant issue is the presence of noninvertible arrows, not the presence of nontrivial automorphisms.

\textbf{Proposition (the canonical comparison).}\ours
For $W=W_\infty$, the ordinary weak homotopy equivalences of categories, the canonical localization
\[
 \rho:\HOT_W(S)\longrightarrow\HOT_W^{\mathrm{lc}}(S)
\]
is an equivalence if and only if $S$ is a groupoid. The same assertion holds for $\sigma$.

We use the established interpretation of the right-fine theory as coherent space-valued presheaves and its locally constant subcategory. Equivalently, one can use Part~V and the split-fibration presentation; the proof below needs only discrete representables for the difficult implication. If $S$ is a groupoid, every functor defined on $S^{\mathrm{op}}$ sends arrows to equivalences. Hence every object is locally constant, and the fully faithful right adjoint $i$ is essentially surjective. Its left adjoint $\rho$ is therefore an equivalence.

Conversely, suppose $\rho$ is an equivalence. For each object $z$, consider the discrete presheaf $h_z(s)=\Hom_S(s,z)$ and its category of elements over $S$. It defines an object of $\HOT_W(S)$. Since this object must be locally constant, precomposition by every $u:a\to b$ induces a weak equivalence of discrete spaces
\[
 \Hom_S(b,z)\xrightarrow{-\circ u}\Hom_S(a,z).
\]
A weak equivalence between discrete spaces is a bijection. With $z=a$, surjectivity gives $v:b\to a$ with $vu=1_a$. With $z=b$, injectivity applied to $uv$ and $1_b$ gives $uv=1_b$, because both become $u$ after precomposition by $u$. Thus every arrow is invertible. Reversing the argument proves the covariant assertion.

For example, take $S=BG$ with $G\ne1$. Its unique object has automorphism group $G$, so $BG$ is not equivalent to a discrete category; nevertheless both canonical comparisons are equivalences. This disproves the printed conjecture even for the ordinary localizer. The proposition concerns the specified comparison functor: detecting a nonconstant representable proves that this functor fails to be an equivalence, without asserting an unrelated classification of possible abstract equivalences between the two categories.

This is our explicit application of Yoneda to the question in the footnote, not a claim of priority. We checked the locally constant comparison in \citet[\S\S1--2]{Cisinski2009LocallyConstant} and the presheaf interpretation in \citet[Theorem~2.2.1.2]{Lurie2009HTT}; neither is being claimed as a new result. We did not find the footnote's proposed criterion stated in the passages checked.
\end{dcomment}

\begin{dcomment}{exposition}{p74:On trouve alors ceci : pour qu'un objet}{Completing the test for local constancy}
The fibre diagram is being used here to recognize objects in an already constructed full subcategory. This is weaker than claiming that an arbitrary diagram in $\Hot_W$ reconstructs the coherent object, an issue treated in the commentary at manuscript p.~7. The recognition statement nevertheless needs the verification requested in footnote~47.

Represent $\xi$ by $X\to S$, and replace it by the split fibration $\Phi X\to S$. Its strict fibre at $s$ is $s\backslash X$, and the transition along $u:s\to t$ is precomposition
\[
 u^*:t\backslash X\longrightarrow s\backslash X.
\]
The replacement map is right-fine, so $X$ and $\Phi X$ represent the same object of $\HOT_W(S)$. If all the $u^*$ belong to $W$, this split fibration belongs to $\underline{\mathrm{Fib}}_0S$: it is a categorical fibration with weakly constant fibres. By the comparison already proved, it belongs to $\underline{\mathrm{Fibhot}}S$. Its class is therefore in the image of $i$. This proves sufficiency without selecting transport maps in an arbitrary nonsplit representative.

For necessity, choose a homotopy-fibration representative $Y\to S$ of an object in the image. Its comma-fibre transition maps are weak equivalences by the defining local-constancy property. The diagram functor sends an isomorphism in $\HOT_W(S)$ to an isomorphism of diagrams in $\Hot_W$. Whether every structural arrow of such a diagram is invertible is invariant under an isomorphism of diagrams. Hence the property holds for every representative of the same class. The covariant argument uses $\Psi X$ and the transitions $X/s\to X/t$.

To pass from ``invertible in $\Hot_W$'' to ``belongs to $W$'', one uses strong saturation. It holds for $W_\infty$ and, in the standard theory, for every basic localizer; see \citet[Proposition~4.2.4]{Cisinski2006Presheaves}. This makes explicit a logical input that the sentence ``cela signifie'' suppresses. Without saturation, the invariant statement uses the saturation of $W$, rather than arbitrary membership in the originally chosen class.

Finally, the factorization through $\Pi S$ in the text concerns only the underlying homotopy-category-valued diagram. Its existence follows from the ordinary universal property of groupoid completion. It does not replace the coherent local system by a representation of the fundamental groupoid. For space-valued coefficients, higher homotopies of the base can still act on the fibres. This leaves the argument above intact: invertibility of all arrows detects local constancy, while the object being detected still retains its coherence.
\end{dcomment}

\begin{dcomment}{examples}{p75}{Maps over a group retain the action}
The common locally constant category still remembers the map to the base. A completely discrete example makes this visible without any higher coherence. Let $G=C_2=\{1,\tau\}$, let $S=BG$, and let $E=G$ with its regular right action. Its category of elements has two objects, with one arrow between any ordered pair; the arrow from $e$ to $e'$ lies over the unique $g$ with $e'=eg$. Its total category is equivalent to $e$, while its projection to $BG$ records the free transitive action.

The associated fibre diagram has one value, the discrete two-point space $E$, with the nontrivial element of $G$ acting by transposition. An endomorphism over the base is an equivariant map $E\to E$. Such a map is determined by the image of $1$, and either choice gives a map: $e\mapsto ae$ for $a\in G$. Thus there are exactly two endomorphisms. Since these diagrams have discrete values, a coherent natural transformation has no nontrivial higher homotopies; their derived mapping space is the same discrete set of two equivariant maps. Consequently the endomorphism set is also two-element in each of the equivalent coarse and fine theories over $BG$.

Forgetting the base would instead replace the total category by a point, whose endomorphism set in $\Ho(\mathrm{Spaces})$ is a singleton. The coarse theory over $BG$ therefore cannot be obtained by simply computing maps between the total homotopy types. It localizes maps \emph{over} $BG$, and the permitted zigzags continue to remember that structure.

The same computation can be made with spaces. The object is the universal covering $EG\to BG$, with contractible total space and discrete fibre $G$. Its derived space of self-maps over $BG$ is $\Omega BG\simeq G$. The ordinary space of self-maps of its contractible total space is contractible. This is the relative distinction behind the Hom formula at manuscript p.~81. It also explains why a groupoid base is not a degenerate case of the theory, even though the fine-to-coarse comparison is an equivalence for such a base.
\end{dcomment}

\begin{dcomment}{exposition}{p75:les 7 carrés du diagramme sont catésiens.}{What the fifteen categories remember}
The expanded diagram separates two choices which are often bundled together in contemporary notation: a property of the projection $X\to S$, and a restriction on the morphisms allowed between such projections. This separation is responsible for both its size and its usefulness.

Except for the categories of split fibrations and split opfibrations, the displayed categories are full subcategories of $\Cat/S$. A map between two fibred objects is therefore an arbitrary functor over $S$; it need not preserve Cartesian arrows. For the split categories, Grothendieck imposes preservation of the splitting, and hence preservation of the chosen Cartesian or cocartesian arrows. Their forgetful functors are not full. The commentary at manuscript p.~30 explains why the naive unit for a splitting construction need not lie in this narrower morphism category.

The seven Cartesian squares assert that the relevant conditions on objects and morphisms are simultaneous conditions. In particular, the square with $\underline{\mathrm{Parf}}S$ expresses the conjunction of the left and right properties, while the outer squares impose the locally constant condition on a split or nonsplit fibration. These assertions should be read in the ordinary category-of-categories sense in which the given subcategories and forgetful functors have been specified. They are not additional assertions that the corresponding nerves form homotopy pullback squares.

The explicit comma replacements remove the possible difficulty caused by the narrow morphism categories. Applied to a functor over $S$, the rule $s\mapsto s\backslash X$ gives a strictly natural transformation of strict $\Cat$-valued diagrams, and therefore a splitting-preserving map of their categories of elements. This supplies a functor even when the original map did not preserve Cartesian arrows. The equivalence after localization follows from the natural comparison zigzags established earlier, not from pretending that all the forgetful functors were already fully faithful. A later systematic treatment of such comma and category-of-elements comparisons is \citet[Theorem~3.8]{Maltsiniotis2005Asphericite}.
\end{dcomment}

\begin{dcomment}{examples}{p78}{A finite comma diagram with an infinite reflection}
The comma replacement and the path replacement can have very different fibres. To see the distinction explicitly, let $S$ have two objects $a,b$ and exactly two nonidentity arrows $u,v:a\rightrightarrows b$. Its nerve realizes to a circle. Take $X=e$ over $b$.

The right comma diagram $F(s)=s\backslash X$ is discrete, with
\[
 F(b)=\{*\},\qquad F(a)=\{u,v\},\qquad
 F(u)(*)=u,\quad F(v)(*)=v.
\]
Its category of elements has three objects, two over $a$ and one over $b$, with one arrow from each object over $a$ to the object over $b$. The latter is terminal, so the total category is contractible. This split fibration is right-fine equivalent to $X$, but it is not locally constant: neither map from a singleton to the two-point set is an equivalence. Dually, $X/s$ is empty at $a$ and a singleton at $b$, since there is no arrow $b\to a$.

We can exhibit a locally constant replacement of $F$ without writing the entire category of paths. Put
\[
 L(a)=L(b)=\mathbb Z,\qquad L(u)(n)=n,
 \qquad L(v)(n)=n+1,
\]
viewing $L$ as a presheaf of discrete spaces. Define $F\to L$ by sending $*$ to $0$, $u$ to $0$, and $v$ to $1$. The two naturality squares commute. All transition maps of $L$ are bijections. Its category of elements has objects $a_n,b_n$ and arrows
\[
 a_n\longrightarrow b_n,\qquad a_{n+1}\longrightarrow b_n
 \qquad(n\in\mathbb Z).
\]
There are no nonidentity composable pairs. Its nerve is the infinite line subdivided into edges, hence is contractible. The functor from the three-object category of elements of $F$ to this infinite category is therefore a coarse weak equivalence. Both its source and target are nonempty and contractible, which verifies the assertion directly rather than by a fibrewise test.

Since $L$ is locally constant, the established reflection identifies its class with $i\rho(F)$. The fibre of this reflection is $\mathbb Z$, whereas the original comma fibres had at most two points. Transport around the generating loop of the circle is translation by one, up to the choice of orientation. Geometrically, the total locally constant object is the universal cover of the circle. The reflection must add all iterates of transport; it cannot be computed by merely declaring the two original transition maps invertible within their finite sets. This example gives a concrete role to the path construction in a place where the comparison diagram alone can conceal it.
\end{dcomment}

\begin{dcomment}{exposition}{p81}{The target controls the comparison of morphisms}
The formula on this page is a universal property of the replacement of the \emph{target}. It is stronger than saying that $X$ and $C(X)$ have the same total homotopy type, and it explains why homotopy fibrations are the right objects in which to receive maps.

Write $Q_{\mathrm d},Q_{\mathrm g},Q$ for the three localization functors out of $\Cat/S$. The adjunction $\rho\dashv i$ identifies
\[
 \Hom_{\HOT_W(S)}(Q_{\mathrm d}X,iQY)
       \cong\Hom_{\HOT_W^{\mathrm{lc}}(S)}(QX,QY).
\]
The chosen model for $iQY$ is $Q_{\mathrm d}CY$, and the analogous model for $jQY$ is $Q_{\mathrm g}CY$. This gives both isomorphisms printed in the box, with their naturality in $X$ and $Y$. One must retain $CY$ on the fine sides unless $Y$ is already local.

For $Y\in\underline{\mathrm{Fibhot}}S$, the map $i_Y:Y\to CY$ is both right-fine and left-fine. Postcomposition by this map identifies the corresponding Hom sets with targets $Y$ and $CY$. Combining these identifications with the box proves footnote~51. In particular, the fine and coarse localizations have the same maps \emph{into a local object}, although they can have different maps between arbitrary objects. This is the familiar defining property of a reflective full subcategory, here recovered from explicit relative-category models.

There is also a useful way to recognize the arrows inverted by $\rho$. A map $f:A\to B$ in $\HOT_W(S)$ is inverted if and only if, for every local object $L$, precomposition induces a bijection
\[
 \Hom(B,L)\longrightarrow\Hom(A,L).
\]
One implication follows from the adjunction. Conversely these bijections say that precomposition by $\rho f$ is a bijection against every object of the coarse category; Yoneda implies that $\rho f$ is invertible. The fibre diagram alone does not supply these Hom sets. The regular-action example above shows why even objects with contractible total spaces can have nontrivial relative endomorphisms.
\end{dcomment}

\begin{dcomment}{exposition}{p82}{Coherence already follows from the two invertibilities}
Grothendieck is trying to decide which part of the reflection depends on the special geometry of paths. Concatenation gives an attractive proof of the equality of the two iterated units, but under the hypotheses already written at manuscript p.~66 that equality is also formal. The distinction matters for the abstraction that follows.

Let $T$ be an endofunctor of an ordinary category and $\eta:1\to T$ a natural transformation. Assume that both $T\eta_X$ and $\eta_{TX}$ are invertible for every $X$. Put $a=T\eta_X$ and $b=\eta_{TX}$. Naturality at $\eta_X$ gives $a\eta_X=b\eta_X$. Applying $T$ and cancelling the invertible map $a$ gives $Ta=Tb$. Next apply naturality to the two maps $a,b:TX\to T^2X$:
\[
 \eta_{T^2X}a=(Ta)b,
 \qquad \eta_{T^2X}b=(Tb)b.
\]
The right sides agree, and $\eta_{T^2X}$ is invertible. Hence $a=b$. Notice that the proof never cancels $\eta_X$, which need not be invertible.

Now work in $\mathcal A\Sigma_0^{-1}$. The functor $C$ induces an endofunctor $T$ because $\Sigma_0\subseteq\Sigma$ and $C(\Sigma)\subseteq\Sigma_0$. Its coaugmentation is the image of $i$. The map $Ti_X$ is invertible because $i_X\in\Sigma$, and $i_{TX}$ is invertible by the second condition at manuscript p.~66. The lemma therefore gives $i_{C(X)}=C(i_X)$ in the fine localization. The first triangle identity is already tautological, and this proves the second. The same argument applies to the left-fine localization.

This is the ordinary-category case of the standard characterization of a localization by an idempotent coaugmented endofunctor; compare \citet[Proposition~5.2.7.4]{Lurie2009HTT}. The displayed proof is included to make the precise cancellation visible. No new geometric property of $C$ is required at this step once both invertibilities are known. Establishing those invertibilities, and identifying the essential image of $C$, remain mathematical work. The source's later Theorem-Scholie~2, at manuscript pp.~108--113, puts exactly these inputs into a general replacement theorem.
\end{dcomment}

\begin{dcomment}{context}{p83}{The extra structure hidden in a categorical base}
The word ``bi-site'' is Grothendieck's way of isolating information that a bare presentation of a homotopy theory does not remember. The categories $S/s$ and $s\backslash S$ retain the orientation of arrows in $S$. Localizing the total category alone can destroy that orientation. Consequently a construction which is invariant under weak homotopy equivalence of the base cannot, without additional input, recover both of the directed diagram theories.


The labels also need a small correction in the surrounding prose. A local
isomorphism is a discrete fibration, modelled by a slice projection $S/s\to S$;
these test $W_S^{\mathrm g}$. A co-local isomorphism is a discrete opfibration,
modelled by $s\backslash S\to S$; these test $W_S^{\mathrm d}$. The sentence at
manuscript p.~82 interchanges the two superscripts. The conventions at pp.~83,
94, and 135 agree with the comma definitions above. The edition's wording is
retained, and the correction register records the swap. This is a notational
slip, not a change in the comparison theorem.

There is a useful historical continuity here. In \citet[\S69, pp.~1--3 of the edited extract]{Grothendieck1983}, Grothendieck describes the passage from one derived category to the family of derived categories of diagrams, relates the motivation to the SGA~5 trace formalism, and then formulates the two Kan-extension adjoints. The present passage asks for a converse kind of explanation: which geometric structure on a category of models permits its objects themselves to serve as indexing bases? For $\Cat$, the answer uses comma categories and the two orientations. It is not supplied merely by knowing that its localization presents spaces.

This also explains the distinction from a slice of a model category. For $S\in\mathcal M$, the ordinary slice model structure, when available, localizes maps by their total weak equivalences. That is the coarse theory in these notes. A category of diagrams indexed by an ordinary category $S$, with objectwise weak equivalences, has a different role: it retains the directed indexing shape. In $\Cat$ the same symbol $S$ can play both roles, which is precisely what makes their comparison possible and initially deceptive.

No Grothendieck topologies or covering axioms for the four adjectives are defined in this paragraph. The expectations about equal cohomology should therefore be recorded as a programme, rather than as a theorem about a specified pair of sites. In particular, the remark identifying ``\'etale'' and ``Zariski'' belongs to these categorical local-isomorphism intuitions; it is not an assertion that the \v{C}ech or sheaf theories of the \emph{scheme-theoretic} \`etale and Zariski topologies coincide. The precise comparison statements available in these pages are the comma and fine-localization equivalences already established.
\end{dcomment}

\begin{dcomment}{development}{p84}{A universal version of the coherence calculation}
For the categorical replacement used here, the question has an affirmative answer under the standing hypotheses of the Theorem-Scholie. Put $U_S=\{u\mid u\times_ST\in W_T\text{ for every }T\to S\}$. Part~III of the Scholie identifies $W_S$ with the fibrewise class $W_S^{\mathrm f}$ on the full subcategory of Hot-fibrant objects. These objects are stable under arbitrary strict base change, as the source states explicitly at manuscript p.~119. If $u:P\to Q$ is a coarse weak equivalence between such objects and $f:T\to S$ is arbitrary, the fibre of its pullback at $t\in T$ is precisely
\[
 (u\times_ST)_t=u_{f(t)}.
\]
It therefore belongs to $W$. The pulled-back source and target remain Hot-fibrant, so Part~III over $T$ gives $u\times_ST\in W_T^{\mathrm f}=W_T$. This proves that $U_S$ and $W_S$ agree between Hot-fibrant objects; the converse uses the identity base change.

The values of $C$ are perfect, hence Hot-fibrant. Naturality of $\eta:1\to C$ and two-out-of-three imply $C(W_S)\subseteq W_S$, and the preceding comparison strengthens this to $C(W_S)\subseteq U_S$. Likewise $\eta_{CX}$ is a coarse weak equivalence between Hot-fibrant objects and therefore lies in $U_S$. Both iterated units are now invertible in $(\Cat/S)U_S^{-1}$, so the cancellation lemma at manuscript p.~82 proves
\[
 C\eta_X=\eta_{CX}\quad\hbox{in }(\Cat/S)U_S^{-1}.
\]
The unit and counit construction below then gives the corresponding reflection. Thus the source's particular question is answered using its fibrewise comparison and base-change stability; no additional identification of strict pullbacks with homotopy pullbacks is needed for that application.

There is also a useful sufficient criterion for other replacements, formulated directly in homotopy-theoretic terms. Here its additional hypothesis is explicit.

\textbf{Proposition (universal comparison for good replacements).}\ours
Fix $S$ and $W=W_\infty$. Let $\mathcal H\subseteq\Cat/S$ be a full subcategory such that, for every $P\to S$ in $\mathcal H$ and every $T\to S$, the canonical square obtained by applying nerves to $P\times_ST$ is a homotopy pullback square. Suppose $C:\Cat/S\to\mathcal H$ has a natural map $\eta:1\to C$ whose components belong to $W_S$. Let $U_S$ consist of those maps $f$ over $S$ for which $f\times_ST$ is a weak equivalence for every $T\to S$. Then
\[
 C(W_S)\subseteq U_S,\qquad
 C\eta_X=\eta_{CX}\quad\hbox{in }(\Cat/S)U_S^{-1}.
\]
Moreover the canonical functor
\[
 (\Cat/S)U_S^{-1}\longrightarrow(\Cat/S)W_S^{-1}
\]
has a fully faithful right adjoint induced by $C$.

To prove the assertion, first consider a coarse weak equivalence $f:P\to Q$ between objects of $\mathcal H$. For any $T\to S$, its base change is represented, up to weak equivalence, by the map between homotopy pullbacks
\[
 NP\mathbin{\times^{h}_{NS}}NT
   \longrightarrow NQ\mathbin{\times^{h}_{NS}}NT.
\]
This is an equivalence because homotopy pullback preserves equivalences of diagrams. Thus $f\in U_S$. The converse holds by taking $T=S$. Consequently $U_S$ and $W_S$ agree on maps between objects of $\mathcal H$.

Naturality and two-out-of-three imply that $C$ preserves $W_S$. Its values lie in $\mathcal H$, so the preceding paragraph gives $C(W_S)\subseteq U_S$. Also $\eta_{CX}$ is a coarse equivalence between two objects of $\mathcal H$, hence belongs to $U_S$. In $(\Cat/S)U_S^{-1}$ the two iterated units are therefore invertible. The cancellation lemma at manuscript p.~82 proves their equality. The unit of the proposed adjunction is the image of $\eta$, and its counit in the coarse localization is the inverse of that image. The two triangle identities follow from this equality and from $\eta^{-1}\eta=1$. Since the counit is invertible, the right adjoint is fully faithful. This proves all assertions.

This second criterion applies only after its homotopy-pullback hypothesis has been established, for example by a relative Theorem~B argument after each base change. It does not infer that hypothesis solely from the word ``fibration'', and it does not assert that every category equipped only with weak equivalences has such a replacement. The later Theorem-Scholie~2, especially manuscript pp.~111--112, supplies the abstract comparison for any finer class agreeing on fibrant objects; the opening argument verifies that agreement for the source's universal class. Our contribution here is the explicit verification for universal equivalences by the two stated methods. We checked the localization criterion in \citet[Proposition~5.2.7.4]{Lurie2009HTT} and the source's later abstract theorem; no priority is claimed for either ingredient.
\end{dcomment}

\begin{dcomment}{status}{p85}{The locally constant inclusion has two adjoints}
The programme has a sharper established form than the reflection proved here. For a derivator $\mathbb D$ and a small category $A$, \citet[Proposition~3.10 and Theorem~3.13]{Cisinski2009LocallyConstant} proves that the full inclusion of locally constant coherent diagrams has both a left and a right adjoint, and that, for fixed $A$, these subcategories form a derivator in an additional diagram variable. This is not the assertion that the varying-base assignment $S\mapsto\HOT_W^{\mathrm{lc}}(S)$ is a derivator. These are the theorem numbers in the published 2009 article; its earlier preprint uses different numbering. The model-category construction is given there by the local structures of \S2, with the fibrant objects characterized in Proposition~2.3.

For spaces, the reason is transparent in the language of homotopy-coherent functors. Let $q:NS\to\Pi_\infty S$ be groupoid completion. Restriction
\[
 q^*:\mathrm{Fun}((\Pi_\infty S)^{\mathrm{op}},\mathrm{Spaces})
       \longrightarrow\mathrm{Fun}((NS)^{\mathrm{op}},\mathrm{Spaces})
\]
is fully faithful and identifies its image with the functors carrying all arrows to equivalences. Its left and right Kan extensions give the two adjoints, since spaces admit small homotopy colimits and limits. Passing to homotopy categories gives the corresponding statement here.

The second adjoint of this one inclusion should not be identified with the functor $j$ in the Scholie. The latter embeds the coarse category into a \emph{different}, covariant fine theory. This keeps the later general theorem distinct from Grothendieck's two parallel reflections. The vague site comparison of the preceding paragraph requires a precise choice of sites before it acquires a similarly definite present status.
\end{dcomment}

\begin{dcomment}{exposition}{p87}{Objectwise replacement and functorial replacement}
The restart of 16 November isolates a familiar but consequential distinction. A weak equivalence $X\to C(X)$ with fibrant target gives a representative of each localized object; it does not by itself define a functor on localized morphisms. Grothendieck correctly stops at essential surjectivity before adding a hypothesis that controls maps.

In the functorial case the proof uses much less than a full model structure. Put $\mathcal A=\mathcal M/S$, let $\mathcal H=\underline{\mathrm{Fib}}S$ be full, and let $D:\mathcal H\hookrightarrow\mathcal A$ be its inclusion. Suppose $C:\mathcal A\to\mathcal H$ is a functor and $\eta:1_{\mathcal A}\to DC$ has components in $W_S$. For $u:X\to Y$, naturality gives
\[
 DC(u)\eta_X=\eta_Yu.
\]
If $u\in W_S$, two-out-of-three implies $C(u)\in W_S$. Hence both functors descend to the localizations. The transformation $\eta$ becomes an isomorphism $1\cong\overline D\,\overline C$. Restricting it to $\mathcal H$ gives $1\cong\overline C\,\overline D$, since fullness ensures that the same components are maps in $\mathcal H$. Thus $\overline C$ and $\overline D$ are quasi-inverse equivalences.

For any full $\mathcal F$ with $\mathcal H\subseteq\mathcal F\subseteq\mathcal A$, the same proof applies to $C|_{\mathcal F}$ and the two inclusions. This proves the intermediate-category assertion at manuscript p.~89. No lifting axiom, finite-limit hypothesis, or base-change stability is used in this functorial branch. Those properties become relevant in the other branch, where lifts have to replace a functorial choice.

The word ``fibrant'' is relative to the chosen base: the objects of $\mathcal H$ are maps $X\to S$ in the distinguished class $F$. It does not assert that $X$ is fibrant as an object of $\mathcal M$, unless $S$ is terminal or an additional argument relates the two conditions. This distinction will also matter when path objects are formed inside the slice in the following pages.
\end{dcomment}

\begin{dcomment}{exposition}{p89}{What makes a comparison canonical}
The desired independence of replacement is an independence with specified comparison maps. This qualification explains the otherwise too compressed phrase ``unique \`a isomorphisme unique pr\`es'' without introducing a new gap in the argument.

Suppose $\delta:\mathcal H\to\mathcal A$ is an equivalence, and $\gamma,\gamma'$ are quasi-inverses equipped with isomorphisms $\epsilon:\delta\gamma\to1$ and $\epsilon':\delta\gamma'\to1$. Full faithfulness of $\delta$ gives a unique natural transformation $\theta:\gamma\to\gamma'$ satisfying
\[
 \delta\theta=(\epsilon')^{-1}\epsilon.
\]
Its components are invertible because their images are, so this is the unique comparison compatible with the specified identifications. This is the precise meaning needed for replacement functors. If the identifications are forgotten, a quasi-inverse can have nontrivial natural automorphisms; for the identity equivalence of the one-object category $BC_2$, its natural automorphisms are the two central elements of $C_2$.

At manuscript p.~90 the chosen map $i_X$ has the \emph{left} lifting property with respect to fibrations in current terminology: it is the left vertical side of the square, while $C(X')\to S$ is the right vertical side. The printed phrase ``rel\`evement \`a droite'' on p.~86 is resolved by this diagram. The lift supplies a possible $\bar u:C(X)\to C(X')$. Its uniqueness has not been asserted. What the next lemma proves is that different lifts give the same localized arrow. Once that is proved, identity and composition follow by choosing the identity lift and the composite lift. This is how a nonfunctorial choice of fibrant objects becomes a genuinely functorial construction after localization.
\end{dcomment}


\begin{dcomment}{context}{p91:Il reste à prouver cette indépendance,}{Choosing replacements without choosing a functor}
Grothendieck is separating the existence of good representatives from the possibility of doing mathematics with them. A factorization of every object over $S$ gives a representative in $\Fib S$, but it does not yet identify the morphisms of the two localizations. That is the work of the next lemma. The question is narrower than the construction of a model category: we require enough lifting to extend a map into a fibrant object, and enough paths to make two extensions agree after localization.

The lifting class is
\[
 \mathrm{Cof}_0={}^{\perp}\mathrm{Fib}.
\]
Thus its maps have the \emph{left} lifting property against fibrations. Neither the notation nor the definition says that every member of this class lies in $W$; this is why the factorization hypothesis subsequently asks for $i_X\in\mathrm{Cof}_0\cap W$. The manuscript's occasional phrase ``relèvement à droite'' describes the wrong side of the lifting square, as the editors explicitly correct on p.~108. For the extension problem here, the right vertical arrow is the structural map $C'\to S$. The square with top arrow $X\to C'$ and bottom arrow $C\to S$ has a lift because $i:X\to C$ belongs to ${}^{\perp}\mathrm{Fib}$. No choice of a cleavage or of a functorial factorization is involved.

What must become canonical is the \emph{class} of the extension, not an actual extension. Once its class is independent of the lift, a composite of two chosen lifts is itself an allowable lift of the composite. Independence therefore supplies compatibility with composition. This is an instructive instance of a construction becoming functorial only after passing to a homotopy category.

In modern language the other route, hypothesis A, is a right deformation: an endofunctor $C$ with a natural weak equivalence $\mathrm{id}\to C$ and values in a full subcategory of good objects. The equivalence of localizations follows from this data, without a model structure; compare the dual of \citet[Definition~2.2.1 and Remark~2.2.2]{Riehl2014}. Hypothesis B replaces this strictly natural data by lifting and a proof of independence. Its advantage is that a model category supplies it even when we do not wish to choose a functorial factorization. Its cost is the relative path-object argument that Grothendieck now gives.
\end{dcomment}

\begin{dcomment}{exposition}{p93}{The path object makes the choices agree}
Grothendieck's proof works inside the subcategory whose localization we want to understand. This point prevents a circular argument: knowing that two lifts become equal in $(\mathcal M/S)[W_S^{-1}]$ would not yet show that they become equal in $(\Fib S)[W_S^{-1}]$, since full faithfulness of the inclusion is precisely what remains to be proved.

Write $\mathcal A=\mathcal M/S$ and $\mathcal F=\Fib S$. Suppose that $v_1,v_2:C\to C'$ agree after precomposition with $i:X\to C$, where $i$ has the left lifting property against fibrations and $C,C'$ lie in $\mathcal F$. Factor the relative diagonal as
\[
 C'\xrightarrow{j}P\xrightarrow{(d_0,d_1)}C'\times_S C',
 \qquad j\in W,\quad (d_0,d_1)\in\mathrm{Fib}.
\]
The object $C'\times_S C'$ is fibrant over $S$: either projection is a pullback of $C'\to S$, and its composite with the other structural map is a fibration. The same is true of $P$. The equalities $d_0j=d_1j=\mathrm{id}_{C'}$ therefore take place in $\mathcal F$. Since $j$ is inverted in $\mathcal F[W^{-1}]$, we obtain $d_0=d_1=j^{-1}$ there.

The square with left side $i$, right side $(d_0,d_1)$, bottom side $(v_1,v_2)$, and top side $jv_1i=jv_2i$ has a lift $h:C\to P$. Consequently $v_1=d_0h=d_1h=v_2$ in $\mathcal F[W^{-1}]$. This is the whole independence argument. In fact, the additional condition $i\in W$ is not needed for this particular conclusion. It is needed when the chosen $i_X$ are used to compare an arbitrary object with its replacement.

For completeness, let $u:X\to Y$ and choose a lift $\bar u:CX\to CY$ satisfying $\bar u i_X=i_Yu$. Define $\gamma(u)=[\bar u]$. The identity of $CX$ is a permissible lift for $\mathrm{id}_X$, and $\bar v\bar u$ is a permissible lift for $vu$. Independence proves the functor laws. If $u\in W$, the equation $\bar u i_X=i_Yu$ and two-out-of-three imply $\bar u\in W$. Thus $\gamma$ descends through the localization at $W_S$.

The same squares finish the naturality checks left on p.~94. In the coarse localization they give a natural isomorphism $\mathrm{id}\xrightarrow{\sim}D\gamma$, where $D$ denotes inclusion. On any full intermediate subcategory $\mathcal F\subseteq\mathcal G\subseteq\mathcal A$, both $X$ and $CX$ belong to $\mathcal G$, so the same maps give the second natural isomorphism $\mathrm{id}\xrightarrow{\sim}\gamma D$. These two isomorphisms establish the comparison theorem without identifying arbitrary fractions by hand. The proof also records exactly where closure of $\mathrm{Fib}$ under pullback and composition is used: in constructing the path object \emph{inside} $\mathcal F$.
\end{dcomment}

\begin{dcomment}{exposition}{p95}{Testing equivalences after change of base}
The class $W_S^{\mathcal L}$ is designed to remember information that a total-space equivalence can lose. It asks that an equivalence survive every observation of the base allowed by $\mathcal L$. At this stage $\mathcal L$ is a collection of tests, not a topology and not a system of weak equivalences. Its role is entirely expressed by pullback.

For a fixed base $S$, the definition can be written
\[
 W_S^{\mathcal L}
 =\bigcap_{(g:T\to S)\in\mathcal L}(g^*)^{-1}(W_T).
\]
Each inverse image contains the isomorphisms and satisfies two-out-of-three; so does their intersection. The identity test gives $W_S^{\mathcal L}\subseteq W_S$. These statements use neither closure of $\mathcal L$ under composition nor closure under base change. The doubt in footnote~57 therefore has a positive answer for this part of the construction.

Composition becomes useful when we move the base. If $g:T\to S$ lies in $\mathcal L$ and $\mathcal L$ is composition-closed, then
\[
 g^*(W_S^{\mathcal L})\subseteq W_T^{\mathcal L}.
\]
Indeed, for a further test $h:U\to T$ in $\mathcal L$, the pullback $h^*g^*u$ is canonically isomorphic to $(gh)^*u$, and $gh$ is an allowed test over $S$. This proves contravariance along the chosen class. Contravariance along an arbitrary $g$ would need the stronger condition that $gh$ be allowed whenever $h$ is allowed; it does not follow from the definition alone.

The localization functor $\rho$ forgets the extra tests. The question about a right adjoint concerns a special class of objects on which those tests no longer distinguish additional weak equivalences. Hypothesis~1 says exactly that: two fibrations over $S$ which are equivalent on total spaces remain equivalent after every allowed base change. It permits the inclusion of fibrant objects to descend to the finer localization. It does \emph{not} say that every map $X\to CX$ is a fine equivalence; that would make the reflection an equivalence and erase the very distinction being introduced.

For the categorical localizers earlier in the chapter, local and colocal tests retain opposite directions of a diagram. Under the standing fundamental-localizer hypotheses, the two suggestions in footnotes~55--56 follow from the base-change theorem discussed at p.~122 below. Every proper base change preserves $W_S^{\mathrm d}$, and hence preserves the total weak equivalence of such a map. Conversely, the coslice projections $s\backslash S\to S$ are proper: they are discrete opfibrations. Testing against all proper maps therefore includes the defining coslice tests and gives exactly $W_S^{\mathrm d}$. The dual argument identifies testing against smooth maps with $W_S^{\mathrm g}$. When properness or smoothness is taken relative to a specified $W$, the corresponding relative base-change criterion is the needed hypothesis. These conclusions do not extend to an arbitrary test class in an arbitrary relative category. The commentary on Chapter~I, manuscript pp.~52--58, and on this chapter, pp.~41--44, explains the geometric criteria; the abstract argument here needs only equality of the two localizers on $\Fib S$.
\end{dcomment}

\begin{dcomment}{exposition}{p97}{Why the properness argument works}
Grothendieck wants an easily checked condition which guarantees that a weak equivalence between fibrations remains an equivalence after a change of base. The lemma gives a sufficient condition in his sparse axiomatic setting. Its two cases are worth keeping distinct, because they apply the same properness axiom to different squares.

Let $u:X\to Y$ be a weak equivalence over $S$, with $X\to S$ and $Y\to S$ fibrations, and let $g:T\to S$. Factor $g$ as
\[
 T\xrightarrow{a}\bar T\xrightarrow{b}S,
 \qquad a\in W,\quad b\in\mathrm{Fib}.
\]
First pull back along $b$. The map $Y\times_S\bar T\to Y$ is a fibration, being a pullback of $b$. Right properness, applied to $u$ along this fibration, says that $u_{\bar T}$ lies in $W$. The structural maps of its source and target over $\bar T$ are still fibrations.

Now pull back along $a$. Right properness applied to those two structural fibrations says that both projections
\[
 X\times_S T\longrightarrow X\times_S\bar T,
 \qquad
 Y\times_S T\longrightarrow Y\times_S\bar T
\]
lie in $W$. In their commutative square the lower horizontal arrow is $u_{\bar T}\in W$. Two-out-of-three proves $u_T\in W$. This completes the argument, with no assumption that $g$ itself be a fibration or a weak equivalence.

The sufficient condition is stronger than necessary in a full Quillen model category. For every map $g:T\to S$, the pullback functor between the slice model categories preserves fibrations and trivial fibrations. It is right adjoint to postcomposition with $g$, and is therefore a right Quillen functor. The dual of Ken Brown's lemma shows that it preserves weak equivalences between objects fibrant in the slice. Thus Hypothesis~1 holds for all base changes in any model category, without global right properness. The base $S$ need not itself be fibrant: fibrancy of $X\to S$ is fibrancy as an object of the slice model category. The corresponding statement for a category of fibrant objects is proved as Lemma~C.1.11 of \citet{RiehlVerity2022Elements}, using Lemma~C.1.10. This does not invalidate Grothendieck's proof: his preceding axioms are weaker than the full model-category axioms, so a separate sufficient condition is useful.

Right properness concerns pulling back an arbitrary weak equivalence \emph{along a fibration}. The conclusion here concerns pulling back a weak equivalence \emph{between two fibrations over the same base} along an arbitrary map. The factorization argument explains the passage from one arrangement to the other.
\end{dcomment}

\begin{dcomment}{exposition}{p100}{Naturality in the finer localization}
The proposed unit of the reflection has to exist where fewer arrows have been inverted. This is the issue hidden by the notation $i_X$: in the coarse localization $i_X$ is invertible, whereas in the finer localization it usually is not. The construction by lifts supplies the required stronger naturality, although the sentence on p.~98 mentions only the coarse category.

Put $\mathcal A=\mathcal M/S$, $\mathcal F=\Fib S$, and write $W_0$ for either $W_S^{\mathcal L}$ or the more general fine localizer of Theorem-Scholie~2. We assume $W_0\subseteq W_S$ and equality on arrows of $\mathcal F$. For a map $u:X\to Y$, a chosen lift satisfies the strict identity
\[
 \bar u i_X=i_Yu
\]
in $\mathcal A$. Two choices of $\bar u$ agree in $\mathcal F[W_S^{-1}]$, hence also after mapping into $\mathcal A[W_0^{-1}]$, because $W_S$ and $W_0$ coincide on $\mathcal F$. The displayed identity therefore defines a natural transformation in the finer localization. Naturality for formal inverses of $W_0$-arrows follows by multiplying a naturality square by those inverses; this is precisely the universal property of localization for natural transformations.

For $u=i_X$, both $\overline{i_X}$ and $i_{CX}$ extend the map $i_{CX}i_X$. The path-object lemma identifies them in $\mathcal F[W_S^{-1}]$. This is stronger than the equality needed in $\mathcal A[W_0^{-1}]$. Let $\eta_X=[i_X]$ be the fine unit, and let the coarse counit be $[i_X]^{-1}$. One triangle identity is cancellation of $[i_X]$ with its inverse. The other is cancellation after the equality
\[
 [C(i_X)]=[i_{CX}]
\]
just proved. Consequently $\rho\dashv i$, and the counit is invertible, so $i$ is fully faithful.

No strict functor $C:\mathcal A\to\mathcal F$ has been manufactured in this argument. The expression $C(i_X)$ means a class represented by an allowable lift. The remark on p.~100 that the entire displayed square lies in $\Fib S$ is a small source slip: its upper-left object $X$ need not be fibrant. The square lies in $\mathcal M/S$; the two arrows whose classes are compared have fibrant sources and targets, which is all the proof uses.
\end{dcomment}

\begin{dcomment}{exposition}{p103}{The mapping-path factorization and its hypotheses}
Grothendieck now wants a replacement with visible geometry. A point of $C_S(X)$ consists of a point of $X$ together with a path from its image to a new point of $S$. The new structural map remembers the endpoint. This changes the map to $S$ while preserving the total homotopy type, and explains why the two projections of the path object play different roles.

Assume finite limits, pullback and composition stability of $\mathrm{Fib}$, and natural maps
\[
 S\xrightarrow{\alpha_S}P(S)
 \xrightarrow{(\sigma_S,\beta_S)}S\times S
\]
with composite the diagonal, $\alpha_S\in W$, and $(\sigma_S,\beta_S)\in\mathrm{Fib}$. Define
\[
 C_S(X)=X\times_{S,\sigma_S}P(S),\qquad
 q:C_S(X)\to X,\qquad
 \bar f=\beta_S\operatorname{pr}_{P(S)}.
\]
The unit sends $x$ to $(x,\alpha_Sf(x))$, so $qi_X=\mathrm{id}_X$ and $\bar f i_X=f$. These formulas describe genuine morphisms: the generalized-element description is justified by the universal property of the pullback and the Yoneda embedding, not by an assumption that the category consists of sets.

To see that $\bar f$ is a fibration, rewrite the pullback as
\[
 C_S(X)=(X\times S)\times_{S\times S}P(S),
\]
where $X\times S\to S\times S$ is $f\times\mathrm{id}$. Its projection $p':C_S(X)\to X\times S$ is a fibration. The second projection $X\times S\to S$ is a pullback of $X\to e$, so it is a fibration if $X$ is fibrant. Hypothesis~3 makes this true for every $X$. Thus $\bar f=\operatorname{pr}_2p'$ is a fibration.

To see that $i_X$ is a weak equivalence, first note that $\sigma_S\alpha_S=\mathrm{id}$ makes $\sigma_S$ a weak equivalence. It is a fibration since $S$ is fibrant and $(\sigma_S,\beta_S)$ is a fibration. Hypothesis~4 then makes its pullback $q$ a trivial fibration. Since $qi_X=\mathrm{id}$, two-out-of-three gives $i_X\in W$. The proof uses pullback stability of \emph{trivial fibrations}; it does not require arbitrary weak equivalences to survive arbitrary pullbacks.

This is the mapping-path mechanism behind Brown's factorization lemma. Its persistence in higher category theory is especially explicit in \citet[Lemma~1.2.19]{RiehlVerity2022Elements}, where a map in an $\infty$-cosmos factors as a section of a trivial fibration followed by an isofibration. The categorical hypotheses and the meaning of equivalence differ, but the pullback calculation is the same.

The concise condition A$'$ on p.~112 must be read as the advertised \emph{strengthening of A}. In particular, it retains the conclusion $i_X\in W$ already required in A. If its path data are instead used as an independent recipe for proving A, the pullback-stability step just exhibited must be supplied, for example by Hypothesis~4. This distinguishes an abbreviated restatement from a missing hypothesis in the theorem as actually formulated.
\end{dcomment}

\begin{dcomment}{examples}{p105}{A groupoid in which both inserted paths are visible}
Grothendieck's element notation is most useful when it allows us to see the two maps $CX\rightrightarrows C^2X$ before localizing. A groupoid gives a completely algebraic example, including a small correction to the order of the formulas on p.~105.

Let $G$ be a group, let $S=BG$ be its one-object groupoid, and let $X=e\to BG$ select the object. Work in groupoids with equivalences as weak equivalences and isofibrations as fibrations. Take $P(S)$ to be the arrow groupoid: its objects are $g\in G$, and a morphism $g\to g'$ is a pair $(a,b)$ with
\[
 bg=g'a.
\]
Source and target send this pair to $a$ and $b$. The constant-path functor sends the unique object of $BG$ to $1\in G$ and an arrow $a$ to $(a,a)$. It is fully faithful, and every $g$ is isomorphic to $1$, so it is an equivalence. The endpoint functor $P(S)\to S\times S$ is an isofibration: prescribed endpoint isomorphisms $(a,b)$ lift by choosing $g'=bga^{-1}$.

The pullback defining $C_S(e)$ fixes the source component of a morphism to be the identity. Its objects are still $g\in G$, but between $g$ and $g'$ there is exactly one morphism, whose image in $BG$ is
\[
 b=g'g^{-1}.
\]
Thus $C_S(e)$ is a contractible groupoid. Its projection to $BG$ is an isofibration: an arrow $h$ with source the image of $g$ lifts to the unique arrow $g\to hg$. The inclusion $e\to C_S(e)$ chooses $1$ and is an equivalence. The strict fibre over the object of $BG$, on the other hand, is the discrete set $G$, since $g'g^{-1}=1$ precisely when $g'=g$. The replacement has preserved the total type while encoding the nontrivial fibre needed over $BG$.

Objects of $C_S^2(e)$ are pairs $(g,h)$, representing successive paths. The two maps are
\[
 i_{C_S(e)}(g)=(g,1),\qquad
 C_S(i_e)(g)=(1,g).
\]
The first inserts a constant path at the \emph{endpoint}; the second first replaces the original object $e$ and therefore inserts a constant path at the \emph{source}. These are the reverse of the labels assigned by ``respectivement'' in the edition. The formulas themselves are present there; their attribution is interchanged. We retain that source slip and use the correctly labelled formulas in the commentary.

Composition gives $\gamma(g,h)=hg$. It is also a functor on morphisms. The unique morphism $(g,h)\to(g',h')$ has image
\[
 d=h'g'g^{-1}h^{-1}
\]
in $BG$, which is exactly the image of the unique morphism $hg\to h'g'$ in $C_S(e)$. Therefore $\gamma i_C=\gamma C(i)=\mathrm{id}$ already before localization. Since the two insertion maps are weak equivalences between fibrant objects, they become the same inverse of $\gamma$ after localization. This example exhibits the geometric comparison which the abstract argument needs; it does not identify categorical fibrations with model-category fibrations in general.

The same example makes a noninvertible fine unit concrete. Let $W_0$ on groupoids over $BG$ consist of equivalences on each strict fibre. It is a two-out-of-three class, obtained by applying the fibre functors. It is contained in the total equivalences for maps between isofibrations, and there it agrees with the total equivalences. To obtain a class contained in the total equivalences on all objects, take its intersection with $W$; the agreement on isofibrations remains true. If $G$ is nontrivial, the unit $e\to C_S(e)$ is a total equivalence but is not an equivalence on the strict fibre: there it is the inclusion of the singleton $\{1\}$ into the discrete groupoid $G$. The fibre functor, followed by localization at equivalences of groupoids, detects that this unit cannot become invertible in the fine localization. Thus the fine and coarse localizations differ even though their equivalences between fibrant objects coincide. This is precisely the situation of the reflection theorem.
\end{dcomment}

\begin{dcomment}{exposition}{p107}{What composition contributes to the reflection}
Grothendieck asks for a composition of paths because it contracts the two-step replacement back to a one-step replacement. For this purpose there is no triple of paths to compare. Associativity therefore has no role in the proof: the only identities used are $c\circ1=c$ and $1\circ c=c$. This explains the deliberately weaker formulation on pp.~106--107.

With the corrected order of the two unit insertions, the formula
\[
 \gamma_X(x,c,c')=(x,c'\circ c)
\]
gives a common left inverse of $i_{CX}$ and $C(i_X)$. Two-out-of-three makes $\gamma_X$ a weak equivalence, because $\gamma_Xi_{CX}=\mathrm{id}$ and $i_{CX}\in W$. Both its source and target are fibrant over $S$, so Hypothesis~1 also makes it a fine weak equivalence. Its inverse in the fine localization is therefore both $i_{CX}$ and $C(i_X)$. The path calculation of manuscript p.~69 has now been isolated from its construction in $\Cat$.

An internal category would require associative composition in addition to the endpoint and unit identities, and its set of arrows on a test object $T$ would be $P(X)(T)$. The edition's $\mathrm{Ch}(T)$ in the paragraph spanning pp.~106--107 is a type error: $\mathrm{Ch}(T)$ is an object of $\mathcal M$, whereas the discussion concerns a set of arrows between elements of $X(T)$. The intended expression is $\mathrm{Ch}(X)(T)$.

There is also a shorter route once the comparison of localizations has been proved. That comparison can transfer the required equality from the coarse localization back to the localized fibrant subcategory. We give the argument at p.~111. It shows that the path composition is a useful geometric witness, while the formal reflection theorem itself needs only the functorial replacement of A. This strengthens the stated theorem without changing the validity of the path proof.
\end{dcomment}

\begin{dcomment}{exposition}{p108}{Reading the two versions of the abstract theorem}
Theorem-Scholie~2 organizes the preceding work around two independent questions. The first is whether fibrant representatives compute the coarse localization. The second is whether those representatives form a reflective subcategory of a finer localization. Keeping the two questions separate makes the scope of each hypothesis transparent.

For I under A, we need a full subcategory $\mathcal F\subseteq\mathcal A$, a functor $C:\mathcal A\to\mathcal F$, and a natural transformation $\mathrm{id}_{\mathcal A}\to C$ whose components are in $W$. Two-out-of-three shows that $C$ preserves $W$. Restricting the transformation to $\mathcal F$ gives the second comparison with the identity, and hence an equivalence of localizations. Finite limits, fibrations, and lifting properties are irrelevant to this formal argument. Their role is to provide the replacement in particular examples.

Under B, the functor only appears after localization. The left lifting property supplies its arrows, and the path-object lemma supplies independence. In the proof as written, finite limits and stability of $\mathrm{Fib}$ under pullback and composition are used at pp.~92--93 to keep the relative path object fibrant over $S$. Footnote~59 says that these stability hypotheses occur only in II, and only in A$'$. That description of the proof is incorrect. It does not affect the theorem, because the hypotheses are included in its statement.

The condition in II is equally precise:
\[
 W_0\subseteq W,
 \qquad W_0\cap\mathrm{Fl}(\mathcal F)
       =W\cap\mathrm{Fl}(\mathcal F).
\]
It is an equality on arrows between good objects, not on all arrows with one good endpoint. In particular, the replacement map $X\to CX$ need not belong to $W_0$. The resulting unit in the fine localization is generally noninvertible, while its image in the coarse localization is invertible. That asymmetry is exactly what permits a nontrivial reflection.

The term ``localement constant'' in the abstract theorem is motivated by the application to $\Cat/S$. For an arbitrary relative category $\mathcal M$, the localization $(\mathcal M/S)[W_S^{-1}]$ has been defined without a notion of transport along arrows of $S$. Its identification with local systems is an additional interpretation supplied by a model of homotopy types. The formal theorem establishes the reflection; it does not by itself define monodromy or a theory of locally constant diagrams.
\end{dcomment}

\begin{dcomment}{later}{p111}{Functorial replacement already gives the reflection}
Grothendieck's construction has supplied more than the final formulation uses. Hypothesis A alone suffices for II; the additional path composition in A$'$ is unnecessary for the adjunction. The missing observation is that I has already proved full faithfulness in the coarse localization. We can use that result to establish coherence without cancelling a noninvertible unit in the fine localization.

Here is the precise argument. Let $\mathcal F\subseteq\mathcal A$ be full, let $W_0\subseteq W$ contain the isomorphisms and satisfy two-out-of-three, and assume that they agree on $\mathcal F$. Suppose given a functor $C:\mathcal A\to\mathcal F$ and a natural transformation $\eta:\mathrm{id}_{\mathcal A}\to C$ with $\eta_X\in W$. Put
\[
 \mathcal D=\mathcal A[W_0^{-1}],\qquad
 \mathcal B=\mathcal A[W^{-1}],\qquad
 \mathcal K=\mathcal F[W^{-1}].
\]
By I, the inclusion $J:\mathcal K\to\mathcal B$ is an equivalence. The inclusion into the fine localization gives $k:\mathcal K\to\mathcal D$. The functor $C$ induces $q:\mathcal D\to\mathcal K$, since it carries $W_0\subseteq W$ into $W$ between objects of $\mathcal F$. The transformation $\eta$ descends to $\mathrm{id}_{\mathcal D}\to kq$.

We first compare the iterated units in $\mathcal K$. Naturality in $\mathcal A$ gives
\[
 C(\eta_X)\eta_X=\eta_{CX}\eta_X.
\]
In $\mathcal B$, $\eta_X$ is invertible, so $C(\eta_X)=\eta_{CX}$. Both arrows are arrows of $\mathcal F$. Since $J$ is faithful, their classes are equal in $\mathcal K$ itself. There is no appeal here to the desired full faithfulness of $k$; the faithful functor is the coarse comparison $J$, whose equivalence was established before $W_0$ entered the argument.

For $Y\in\mathcal F$, the arrow $\eta_Y:Y\to CY$ is a weak equivalence in $\mathcal F$. Define a natural counit $\varepsilon:qk\to\mathrm{id}_{\mathcal K}$ by $\varepsilon_Y=[\eta_Y]^{-1}$. Naturality follows by inverting the natural isomorphism induced by $\eta|_{\mathcal F}$. The triangle identities are now explicit:
\[
 k(\varepsilon_Y)\eta_{kY}=\mathrm{id}_{kY},
 \qquad
 \varepsilon_{qX}\,q(\eta_X)
   =[\eta_{CX}]^{-1}[C(\eta_X)]
   =\mathrm{id}_{qX}.
\]
Thus $q\dashv k$. Its counit is an isomorphism, so $k$ is fully faithful. Composing on the left with the equivalence $J$ identifies $q$ with the canonical localization $\rho:\mathcal D\to\mathcal B$ and gives
\[
 \rho\dashv kJ^{-1}.
\]
The functor $kJ^{-1}$ is Grothendieck's $i$. In particular,
\[
 \Hom_{\mathcal D}(X,CY)
 \simeq \Hom_{\mathcal B}(X,Y)
\]
under A without A$'$, and for every choice of $W_0$ with the stated agreement property.

The bijection can also be written without invoking the adjunction as a black box. For $Y\in\mathcal K$, it sends $f:X\to kY$ to $\varepsilon_Yq(f):qX\to Y$. Its inverse sends $h:qX\to Y$ to $k(h)\eta_X$. Substituting one formula into the other, naturality of $\eta$ or $\varepsilon$ moves the comparison map past $f$ or $h$, and the corresponding triangle identity cancels the remaining pair. These formulas establish both naturality variables and show that the comparison with the coarse localization is the canonical one. In particular, no new choice of a map from $X$ into a fibrant target has been made.

This strengthening belongs to the general theory of derived adjunctions, rather than being a new phenomenon. Apply the deformable-adjunction theorem, \citet[Theorem~2.2.11]{Riehl2014}, to the identity adjunction between $(\mathcal A,W_0)$ and $(\mathcal A,W)$: its left functor is already homotopical, and its right functor admits the right deformation $C$, on whose image the two classes agree. The proof above is included to make the application independent of a stronger axiomatization of homotopical categories, and to show precisely how the reflection lemma of manuscript p.~67 receives its coherence condition here.
\end{dcomment}

\begin{dcomment}{examples}{p113}{Different fine localizers can have the same fibrant targets}
Grothendieck's corollary says more than full faithfulness on pairs of fibrant objects. The source $X$ is arbitrary: it is only the \emph{target} $Y$ that must be fibrant. The canonical map of Hom sets is bijective because $Y\to CY$ is already inverted in the finer localization, and the reflection computes maps into $CY$. This is why the Hom set in footnote~60 is independent of the chosen fine localizer.

A finite-set example isolates the statement from homotopical geometry. Let $\mathcal A$ be a skeleton of finite sets, let $W$ consist of all maps, let $\mathcal F$ consist only of a singleton $1$, and put $C(X)=1$. The unique maps $X\to1$ form a natural transformation in $W$. The coarse localization is equivalent to the terminal category: after all arrows are inverted, $X\to1$ is a natural isomorphism to the constant functor. The full subcategory $\mathcal F$ also localizes to the terminal category.

First take $W_0$ to be the bijections. Its localization is just the category of finite sets. The reflection to the terminal category has right adjoint selecting $1$, since
\[
 \Hom_{\mathcal A}(X,1)=\{*\}
\]
for every finite $X$, including the empty set. It plainly does not follow that the fine localization is terminal: for example $\Hom(1,\{0,1\})$ has two elements.

For a second fine class, let $W_1$ consist of the maps which preserve whether a finite set is empty. Equivalently, $W_1$ contains the identity of the empty set and every map between nonempty finite sets. It is the inverse image of the isomorphisms under the functor
\[
 E:\mathcal A\to[1],\qquad
 E(\varnothing)=0,\quad E(X)=1\quad(X\ne\varnothing),
\]
so it satisfies two-out-of-three. The localization at $W_1$ is equivalent to $[1]$. Indeed, the functor $j:[1]\to\mathcal A$ selecting $\varnothing$ and $1$ satisfies $Ej=\mathrm{id}$, and the natural maps $X\to jE(X)$, identity on the empty set and the unique map to $1$ otherwise, all lie in $W_1$. These exhibit quasi-inverse functors after localization.

Thus the two fine categories are $\mathrm{FinSet}$ and $[1]$, which are very different. Nevertheless, both give a singleton set of maps from any object into the distinguished fibrant target. Both $W_0$ and $W_1$ agree with $W$ on $\mathcal F$, so the corollary predicts exactly this agreement. The example also makes the direction of the result unmistakable: maps \emph{out of} the distinguished object are not fixed. In $\mathrm{FinSet}$ there are two maps $1\to\{0,1\}$; in $[1]$ their images coincide.
\end{dcomment}

\begin{dcomment}{later}{p114}{The model-category route in Cat}
Grothendieck invokes Thomason to explain that the abstract replacement argument is not restricted to topological or simplicial objects. Thomason's model structure has as weak equivalences the functors whose nerves are weak homotopy equivalences, and as fibrations the functors detected as Kan fibrations by $\mathrm{Ex}^2N$; its existence is Theorem~4.9 of \citet{Thomason1980}. Factorization into a trivial cofibration followed by a fibration supplies hypothesis B. The theorem concerns the ordinary homotopy theory of nerves; it should not be cited as constructing a model structure for an unspecified class $W$.

There is a historical correction relevant to the preceding properness discussion. Thomason's original proof asserted that Dwyer maps are closed under retracts and hence that every cofibration is a Dwyer map. \citet[pp.~228--229]{Cisinski1999Dwyer} gives a counterexample. His Lemma~4 and Proposition~2 on p.~231 replace Dwyer maps by pseudo-Dwyer maps and recover properness. The model-category existence theorem is unaffected. Thus the later correction repairs an auxiliary assertion in the cited literature, not the validity of Grothendieck's abstract comparison theorem.

The two uses of ``fibration'' in this chapter must remain distinct. A Grothendieck fibration over $S$ carries cartesian lifts of arrows and presents a contravariant diagram of categories. A Thomason fibration is defined by a lifting condition after applying $\mathrm{Ex}^2N$. A categorical comparison between their localized presentations is a theorem, not a consequence of the shared terminology. This is the reason the eight outer categories in the large comparison diagram do not disappear merely because one has a model structure on $\Cat$: they encode strict or split diagrammatic data which the abstract factorization axioms have not named.

The return to those categories is therefore mathematically necessary. The abstract detour has made the localizations and reflection intelligible; the split-fibration presentation will now make variation in $S$ intelligible. These are complementary uses of the same models, and neither replaces the other.
\end{dcomment}

\begin{dcomment}{exposition}{p116}{Pullback is precomposition in the split presentation}
Grothendieck returns to split fibrations because a change of base is completely visible in that presentation. For a strict functor $F:S^{\mathrm{op}}\to\Cat$, write $\int_SF$ for its Grothendieck construction. An object is a pair $(s,x)$ with $x\in F(s)$; an arrow $(s,x)\to(t,y)$ consists of $a:s\to t$ and a map $x\to F(a)y$ in $F(s)$. The projection to $S$ is the associated split fibration.

If $f:S'\to S$, there is a canonical isomorphism over $S'$
\[
 S'\times_S\int_SF\ \cong\ \int_{S'}(Ff^{\mathrm{op}}).
\]
On objects it sends $(s',x)$ with $x\in F(f(s'))$ to the same pair. On an arrow $a:s'\to t'$ it uses the same map $x\to F(f(a))y$. Identities and composition agree because $F$ and $f$ are functors. The chosen cartesian lifts are carried to the chosen cartesian lifts, so the isomorphism also respects the splittings.

An objectwise weak equivalence of functors remains objectwise a weak equivalence after precomposition: its component at $s'$ is the original component at $f(s')$. This proves that the operation descends to the relevant localization and gives the formula for $f^*$ on split representatives. If $g:S''\to S'$ is another functor, precomposition satisfies
\[
 (Ff^{\mathrm{op}})g^{\mathrm{op}}
   =F(fg)^{\mathrm{op}}
\]
strictly. After changing to equivalent presentations by arbitrary fibrations or by categories over $S$, these equalities are expressed by the canonical comparison isomorphisms for iterated pullbacks. The coherence comes from the same strict composition law, rather than from a new choice of replacement for each composite.

This paragraph only establishes the formula on the specified good representatives. Extending it to other objects or deciding when an arbitrary pullback computes the derived inverse image requires the base-change results on the following pages. The categorical calculation here explains why Grothendieck keeps the split models in the theory even after showing that their localized category is independent of the splitting.
\end{dcomment}


\begin{dcomment}{exposition}{p117}{Restriction must be calculated on a suitable model}
Grothendieck now makes the base category vary. The comparison theorem has done more than identify several presentations of one category: it provides presentations on which an operation can be performed before localization. For inverse image the useful presentation consists of fibrations, or more generally smooth categories over the base. The calculation then reduces to the familiar operation on fibres.

Write $q_S:\Cat/S\to\HOT_W(S)$ for localization at $W_S^{\mathrm d}$. If $X\to S$ is smooth, the strict fibre $X_s$ computes the same homotopy type as the comma category $s\backslash X$. For $f:S'\to S$ and $s'\in S'$, the strict fibre of $X\times_S S'$ at $s'$ is exactly $X_{f(s')}$. Consequently a map between smooth objects which is a $W$-equivalence on comma categories remains so after this pullback. The comparison with split fibrations identifies the resulting localized operation with restriction of category-valued presheaves. This proves Proposition~1, including its naturality in maps between smooth objects.

The qualification on the objects matters even when the morphism being pulled back already belongs to $W_S^{\mathrm d}$. Its membership says something about all comma fibres; it says nothing by itself about its strict fibres. Pullback along the inclusion of an object extracts a strict fibre. Thus arbitrary pullback cannot act directly on the entire relative category $(\Cat/S,W_S^{\mathrm d})$. The commentary at manuscript p.~42 (``When strict fibres compute a comma fibre'') isolates precisely the comparison required here.

There is also a useful distinction between the arrows allowed in this presentation and the arrows needed to describe transport. For calculating the category $\HOT_W(S)$, all functors over $S$ between fibrations are allowed; Cartesian functors form a smaller presentation which gives the same localization. A chosen cleavage is therefore not part of a relative homotopy type. When a natural transformation between changes of base is considered on ms.~128--130, a cleavage will nevertheless provide a convenient formula for the induced map. Independence of that auxiliary choice is a separate assertion, not an additional restriction on the morphisms of $\HOT_W(S)$.
\end{dcomment}

\begin{dcomment}{status}{p120:On a le choix entre $\Phi$ et $C$}{The two replacements do different work}
Grothendieck wants one formula for the inverse image of an arbitrary category over $S$. The comma replacement $\Phi_S$ supplies it. The path replacement $C_S$ supplies a formula for a different object, namely the reflection into locally constant relative types. At this point the text briefly identifies those two tasks. Both the sentence preceding Corollary~2 and its second displayed isomorphism are present in the edition; the discrepancy is not introduced by the transcription.

The distinction is already established in the source. The natural map
\[
 X\longrightarrow\Phi_S(X)
\]
belongs to $W_S^{\mathrm d}$, whereas $X\longrightarrow C_S(X)$ is generally only in $W_S$. Accordingly the valid general formula is
\[
 f^*q_S(X)\simeq q_{S'}\bigl(\Phi_S(X)\times_S S'\bigr).
\]
If $i_S:\HOT_W^{\mathrm{lc}}(S)\to\HOT_W(S)$ denotes the inclusion, the formula involving $C_S$ is
\[
 f^*i_S\rho_S q_S(X)
 \simeq q_{S'}\bigl(C_S(X)\times_S S'\bigr).
\]
Here restriction preserves locally constant objects, so that $f^*i_S$ identifies with $i_{S'}f^*_{\mathrm{lc}}$. These equations become the same only when $q_S(X)$ is already locally constant.

A small calculation makes the error decisive without repeating the examples given in the commentary at manuscript pp.~8 and~32. Take $W=W_\infty$, $S=(0\to1)$ and $X=\{0\}\to S$. The comma fibres of $X$ are a point at $0$ and the empty category at $1$. Its relative type is therefore the presheaf of spaces $(\varnothing\to *)$. The local-constant reflection is the constant point: maps from this presheaf to a locally constant arrow are determined by the value at $0$, and evaluation at $0$ identifies locally constant arrows with spaces. Pullback to $1$ consequently gives $\varnothing$ from $X$ and $*$ from $C_S(X)$. Corollary~2 cannot identify them.

This is a localized mistake in the calculation, not a defect in the comparison theorem. The later discussion on ms.~123--127 correctly distinguishes $\rho_Sq_S(X)$ from $q_S(X)$, and the construction of the adjunction on ms.~139--145 consistently uses $\Phi_S$. Keeping the printed error visible is useful: it shows exactly why the two replacement functors were needed.
\end{dcomment}

\begin{dcomment}{exposition}{p121}{A homological test for colocal equivalence}
Grothendieck introduces relative homology here because it packages all the comma-category tests into a single morphism of coherent diagrams. Let $p_X:X\to S$, $p_Y:Y\to S$, and let $u:X\to Y$ be a functor over $S$. For $a\in\mathbf D(e)$, write $a_X$ for the constant coefficient. The map of interest is
\[
 (p_X)_!a_X\longrightarrow(p_Y)_!a_Y.
\]
Its component at $s$ is the map on homotopy colimits of the constant coefficient over $s\backslash X$ and $s\backslash Y$. This is the pointwise Kan-extension formula in the contravariant convention used throughout these pages. If isomorphisms are detected by evaluation, requiring the displayed map to be invertible for every $a$ is exactly the condition that all these comma-category maps be $\mathbf D$-equivalences.

Now suppose $g:T\to S$ is $\mathbf D$-proper in the sense that the relevant base-change comparisons for homological direct image are invertible. Pulling back the displayed map gives
\[
 (p_{X_T})_!a_{X_T}\longrightarrow(p_{Y_T})_!a_{Y_T},
 \qquad X_T=X\times_S T,\quad Y_T=Y\times_S T.
\]
The first map was invertible, and restriction preserves isomorphisms; hence the new map is invertible. Evaluating at every object of $T$ gives the claimed stability of colocal equivalence. This proves the lemma without assuming that strict fibres compute comma fibres: the base-change theorem supplies that comparison at the level of relative homology.

The four classes of test maps listed in the text are genuinely nested. A representable coslice projection is a discrete opfibration, hence a co-local isomorphism; those maps are proper, and absolute properness implies $\mathbf D$-properness. Testing against the largest class therefore implies testing against the first. Conversely the lemma shows that the first condition survives the largest class, and a colocal equivalence is in particular a global $\mathbf D$-equivalence. The source's phrase ``three types'' precedes four entries and has no mathematical consequence.

The historical terminology keeps the geometrical motivation visible. Properness and smoothness were designed to distinguish the two positions in a base-change square, following the analogous cohomological theorems in algebraic geometry. For the categorical definitions, including their dual formulations by taking mates, see \citet[Definition~1.1]{Maltsiniotis2005Asphericite}. Here the lesson is particularly concrete: the general derivator theorem predicts the correct class of changes of base, but cannot yet be used to establish the axioms of the very derivator under construction.
\end{dcomment}

\begin{dcomment}{later}{p122}{The proper base-change conjecture has an elementary proof}
The conjecture on ms.~122 is not an open question. Grothendieck himself proves it later in Chapter~VI, ms.~168--169, printed p.~103 of the complete edition; the footnote points to that proof. The argument avoids the circular appeal to a derivator whose axioms are still being established \citep[ms.~168--169]{Grothendieck1990VI}.

For clarity, isolate its mechanism. If $v:T\to U$ is a universal $W$-equivalence and $a:A\to B$ is a map over $U$, there is a pullback square
\[
\begin{tikzcd}
A\times_U T\arrow[r,"a_T"]\arrow[d] & B\times_U T\arrow[d]\\
A\arrow[r,"a"] & B.
\end{tikzcd}
\]
Both vertical maps belong to $W$. The two-out-of-three property consequently gives $a_T\in W$ if and only if $a\in W$.

Apply this with $g:S'\to S$ proper and an object $s'\in S'$, putting $s=g(s')$. The properness criterion says that
\[
 s'\backslash S'\longrightarrow s\backslash S
\]
is a universal homotopy equivalence; in the relative variant it is a universal $W$-equivalence. A map in $W_S^{\mathrm d}$ is a $W$-equivalence after pullback to $s\backslash S$. The preceding square shows that it remains so after further pullback to $s'\backslash S'$. This is exactly the definition of its pullback belonging to $W_{S'}^{\mathrm d}$. Passing to opposite categories gives the smooth, left-fine version.

Corollary~3 therefore gives the expected strict formula for proper changes of base. Its proof does not depend on choosing a preferred representative: the functor defined on all of $\Cat/S$ agrees, on the smooth presentation, with the previously constructed restriction functor. The equivalence of the two presentations then identifies their descents. The later proof establishes the sufficient condition. The converse requires the additional test calculation given below, and its formulation depends on the chosen localizer.
\end{dcomment}

\begin{dcomment}{exposition}{p126}{Saturation supplies the missing converse}
Grothendieck wants to recognize when restriction commutes with the reflection into locally constant objects. His forward implication is sound: preservation of coarse weak equivalences on any of the indicated smooth or fibred presentations lets the square of localizations descend. For the converse he notices precisely what has not yet been justified. An arrow becoming invertible after localization need not, for an arbitrary chosen class of arrows, have belonged to that class.

There is a short repair once $W$ is strongly saturated in $\Cat$. Let $\mathcal F$ be any category of categories over a fixed base, with any chosen class of allowed morphisms, and let $W_{\mathcal F}$ consist of those allowed morphisms whose underlying functors belong to $W$. The forgetful functor descends to
\[
 \mathcal F[W_{\mathcal F}^{-1}]\longrightarrow\Cat[W^{-1}].
\]
If an actual arrow $u$ of $\mathcal F$ becomes invertible on the left, its underlying functor becomes invertible on the right. Strong saturation gives $u\in W$. The converse is part of localization. Thus the restriction of a strongly saturated coarse class remains strongly saturated, even when $\mathcal F$ is not full.

Assume now that the square in condition~a) commutes. For $u\in W_S$ in a smooth or fibred presentation, its image under $\rho_S$ is an isomorphism. Commutativity says that $f^*u$ becomes an isomorphism in the coarse localization over $S'$. The preceding argument then gives $f^*u\in W_{S'}$. This proves each of b), b$'$), and b$''$). For $W_\infty$, the needed saturation follows immediately by applying the nerve and the homotopy category of simplicial sets. In the later theory it holds for every fundamental localizer; see \citet[Proposition~4.2.4 in the author PDF]{Cisinski2006Presheaves}. The gap belongs to the proof available at this point, not to the resulting equivalence.

Commutation itself is a real restriction. In the interval calculation above, let $f:\{1\}\to(0\to1)$ and let $F=(\varnothing\to*)$. Then
\[
\rho_{\{1\}}f^*F=\varnothing,
 \qquad f^*_{\mathrm{lc}}\rho_S F=*.
\]
This supplies a concrete obstruction without confusing it with the separate failure of arbitrary strict pullback to preserve $W_S^{\mathrm d}$. Grothendieck's Hot-fibration hypothesis is sufficient because it ensures preservation of the coarse class by pullback. Strong saturation closes his converse in Proposition~2; it does not assert that every functor satisfies its equivalent conditions.
\end{dcomment}

\begin{dcomment}{development}{p126}{Properness and a three-object counterexample}
The question on ms.~123 and the proposed counterexample on ms.~126 can both be settled concretely. The useful characterization of $W$-properness is the dual of \citet[Theorem~1.2(b)]{Maltsiniotis2005Asphericite}. For $f:A\to S$, $a\in A$, and $h:f(a)\to b$ in $S$, let $L_f(a,h)$ have objects arrows $k:a\to a'$ with $f(a')=b$ and $f(k)=h$. A morphism is an arrow $\ell:a'\to a''$ with $f(\ell)=\id_b$ and $\ell k=k'$. Then $f$ is $W$-proper exactly when every $L_f(a,h)$ is $W$-aspherical. These are the outgoing comma categories for the inclusion $A_b\to A/b$.

\textbf{Proposition (detecting properness by inverse image).}\ours
For $W=W_\infty$, strict pullback along a fixed $f:A\to S$ preserves right-fine weak equivalences on all categories over $S$ if and only if $f$ is proper. For an arbitrary fundamental localizer $W$, $W$-properness is equivalent to this preservation property holding after every change of base of $f$.

\textit{Proof.} Properness implies preservation by the preceding block, and properness is stable under base change. To prove the first converse, test the assumed preservation on the unit
\[
 \{b\}\longrightarrow\Phi_S(\{b\})=S/b.
\]
This is right-fine: its comma category at $s$ maps the discrete set $\Hom_S(s,b)$ into the category of factorizations of those arrows, and each factorization category has an initial and a terminal object.

After pullback along $f$, take the comma category at $a$. In the target, objects are triples
\[
 (a',k:a\to a',r:f(a')\to b).
\]
The composite $r f(k)$ defines a decomposition indexed by $h\in\Hom_S(f(a),b)$. The component indexed by $h$ has initial object $(a,\id_a,h)$. The corresponding source component is exactly $L_f(a,h)$. A weak equivalence of these disjoint unions to the contractible target components forces every $L_f(a,h)$ to be weakly contractible. The cited criterion gives properness.

For the general converse, suppose preservation holds after every base change. Given $h:f(a)\to b$, pull back $f$ along the functor $[1]\to S$ representing $h$. Apply the preservation hypothesis to $\{1\}\to[1]$ over $[1]$, which is the same unit test. At the object $(a,0)$ of $A\times_S[1]$, its source comma category is $L_f(a,h)$ and its target comma category has an initial object. Since the map belongs to $W$, two-out-of-three implies $L_f(a,h)\to e\in W$. The criterion again gives $W$-properness.\hfill$\square$

\textbf{Example (why arbitrary localizers require base change).}\ours
The distinction between the two formulations cannot be omitted for arbitrary fundamental localizers. Let $W_{\mathrm{ne}}$ consist of functors whose source and target are either both empty or both nonempty. This is a fundamental localizer: weak saturation and the terminal-object axiom are immediate, and the relative theorem-A axiom follows because an object of a nonempty target gives a nonempty target comma category, hence a nonempty source comma category. Take a nontrivial group $G$ and $f:e\to BG$. Over the one-object category $BG$, every comma category of $X\to BG$ is nonempty exactly when $X$ is nonempty. Right-fine $W_{\mathrm{ne}}$-equivalence therefore detects only nonemptiness. Strict pullback along $f$ does not change the object set of $X$, so it preserves that class. Yet $L_f(*,g)$ is empty for every nonidentity $g\in G$, and $f$ is not $W_{\mathrm{ne}}$-proper. This example concerns the broader localizer question, not a weakening of the stronger hypotheses used elsewhere in this chapter.

\textbf{Example (a Hot-fibration which is not proper).}\ours
Let $E$ be the three-object poset with nonidentity arrows $u\to v$ and $u\to w$, and put
\[
 p:E\longrightarrow(0\to1),\qquad p(u)=p(v)=0,\quad p(w)=1.
\]
This is a Grothendieck fibration. The lift $u\to w$ of $0\to1$ is Cartesian: for an object over $0$, its arrows to $u$ and its arrows to $w$ are both singletons when the object is $u$, and both empty when it is $v$. The two fibres are $(u\to v)$ and $\{w\}$; their transition functor chooses $u$ and is a weak equivalence. Thus $p$ lies in $\Fib_0$, hence in $\Fibhot$ under the standing comparison hypotheses.

There is also a direct verification of its particularly strong homotopical behaviour. The section choosing $u$ over $0$ and $w$ over $1$ has a natural transformation $sp\Rightarrow\id_E$ over the base, whose only nonidentity component is $u\to v$. This deformation retraction survives every strict base change. Every projection $E\times_S T\to T$ is therefore a nerve weak equivalence. Two-out-of-three in a square over $S$ shows that pullback along $p$ preserves all coarse weak equivalences.

Nevertheless $p$ is not proper, since $L_p(v,0\to1)$ is empty. More explicitly, the right-fine weak equivalence $\{1\}\to[1]$ pulls back to $\{w\}\to E$. Its comma functor at $v$ is $\varnothing\to e$, so it is not right-fine. This is precisely the counterexample requested on ms.~126. We checked the properness criterion just cited; the elementary tests and examples here are explicit completions of the manuscript's questions, without a priority claim.
\end{dcomment}

\begin{dcomment}{development}{p128}{Local constancy and the missing comma formula}
Grothendieck asks whether the locally constant part, as the base varies, is itself a derivator. His expectation that it is not is correct for ordinary homotopy types, but the suggested obstruction to the existence of adjoints is not the obstruction. It is the pointwise formula for those adjoints that fails. We give the small calculation because it separates this question from the comparison theorem.

\textbf{Proposition (the interval obstruction).}\ours
Let $\mathcal L(S)$ be the full subcategory of the space derivator on $S$ consisting of locally constant objects, with the induced restriction functors. Then $S\mapsto\mathcal L(S)$ does not satisfy the pointwise left Kan-extension axiom.

\textit{Proof.} We use the source's presheaf convention. Let $S=(0\to1)$, let $a:\{0\}\to S$, and let $b:\{1\}\to S$. An object of $\mathcal L(S)$ is a coherent arrow $F_1\to F_0$ which is an equivalence. Evaluation at $0$ is an equivalence
\[
 a^*:\mathcal L(S)\longrightarrow\mathcal L(e)=\Hot.
\]
Indeed the constant-arrow functor is a quasi-inverse: the diagram map from $F_1\to F_0$ to $F_0\xrightarrow{\id}F_0$ has components $F_1\to F_0$ and $\id_{F_0}$, both weak equivalences. This construction is natural in $F$, before taking the homotopy category. Thus a left adjoint $a^{\mathcal L}_!$ exists and is the constant-arrow functor, up to the unique isomorphism between adjoints.

Take the one-point space $*$. We obtain
\[
 b^*a^{\mathcal L}_!(*)\simeq *.
\]
On the other hand, the comma category $1\backslash\{0\}$ of arrows $1\to a(0)$ is empty. The pointwise formula for a left Kan extension in the presheaf convention would therefore identify this value with the homotopy colimit of the empty diagram, namely $\varnothing$. The categories $\mathcal L(e)=\Hot$ and $\mathcal L(\varnothing)=e$ have the usual initial-object adjunction, so that this last value is forced. Since $\varnothing$ is not isomorphic to $*$, the required base-change map cannot be invertible.\hfill$\square$

This proof also identifies what the unrestricted direct image does: in the full space derivator $a_!(*)$ is the representable presheaf $S(-,0)$, with values $*$ at $0$ and $\varnothing$ at $1$. It is not locally constant. The locally constant left adjoint first takes this extension and then reflects it. The reflection changes the value at $1$, and the comma formula is lost.

\textbf{Corollary (commutation beyond Hot-fibrations).}\ours
The inclusion $a:\{0\}\to(0\to1)$ commutes with local-constant reflection for spaces, but is not a Hot-fibration.

\textit{Proof.} The reflection of any arrow $F_1\to F_0$ is the constant arrow on $F_0$. To see its universal property, map to a constant arrow on $K$; the space of maps is $\operatorname{Map}(\operatorname{hocolim}_{S^{\mathrm{op}}}F,K)$, and $0$ is terminal in $S^{\mathrm{op}}$, so this homotopy colimit is $F_0$. Evaluation $a^*$ at $0$ therefore commutes with reflection. Nevertheless the weak equivalence $\{1\}\to S$ pulls back along $a$ to $\varnothing\to\{0\}$, which is not a weak equivalence. Hence $a$ is not a Hot-fibration. This answers the question on ms.~127.\hfill$\square$

The later general result is stronger than is needed for the example. \citet[Corollary~3.11]{Cisinski2009LocallyConstant} proves that restriction on locally constant coefficients has both adjoints for every derivator. His Theorem~3.13 concerns a different construction: one fixes a category $A$ and varies an additional diagram category $I$, retaining arbitrary coherent $I$-diagrams of locally constant $A$-coefficients. That construction is a derivator. It does not say that the varying-base assignment $S\mapsto\mathcal L(S)$ is one.

We checked those statements, the pointwise formula in \citet[Proposition~1.26]{Groth2013}, and the localization description in Cisinski's paper. We did not find this exact interval calculation stated there. The marked result records our explicit completion of the question in these pages; it makes no claim of priority for the obstruction.
\end{dcomment}

\begin{dcomment}{exposition}{p130}{A natural transformation is realized by transport}
Grothendieck has explained restriction on objects; a prederivator also remembers how restriction varies with a natural transformation of indexing functors. For $u:f\Rightarrow g:S'\to S$, the presheaf convention requires a transformation $u^*:g^*\Rightarrow f^*$. A pullback square alone supplies no arrow in this direction. A fibration supplies it through Cartesian lifting.

Let $p:X\to S$ be a Grothendieck fibration. For every $a\in S'$ and $x\in X_{g(a)}$, choose a Cartesian arrow
\[
 \lambda_{a,x}:u_a^*x\longrightarrow x
\]
over $u_a:f(a)\to g(a)$. Send $(a,x)$ in $g^*X$ to $(a,u_a^*x)$ in $f^*X$. For an arrow $(h,v):(a,x)\to(b,y)$ in $g^*X$, the equality
\[
 g(h)u_a=u_bf(h)
\]
and the Cartesian property of $\lambda_{b,y}$ give a unique arrow $w:u_a^*x\to u_b^*y$ over $f(h)$ satisfying
\[
 \lambda_{b,y}w=v\lambda_{a,x}.
\]
Uniqueness proves preservation of identities and composition. This constructs the requested functor over $S'$.

Different choices of Cartesian lifts are related by unique vertical isomorphisms. The resulting transport functors are therefore naturally isomorphic. The same uniqueness gives the comparison between transport along a composite transformation and successive transport, together with its coherence. A splitting can make selected comparison maps identities, but the localized operation is independent of the splitting.

For an arbitrary functor $H:X\to Y$ over $S$, transport need not commute strictly with $H$: $H$ need not preserve Cartesian arrows. Factoring $H(\lambda_{a,x})$ through a chosen Cartesian lift in $Y$ gives a vertical natural transformation between the two resulting functors. A vertical natural transformation is a homotopy over $S'$: it is represented by a functor from the cylinder $X\times[1]$, whose projection to $X$ is a weak equivalence on fibres. Thus the two maps agree after localization. This supplies the naturality needed when all functors over the base, rather than just Cartesian functors, are used in the fibred presentation.

The source's negative observation is accordingly about finding a direct functor on insufficiently structured representatives. It does not obstruct the existence of the $2$-functor. That $2$-functor is already defined by precomposition on coherent presheaves; the lifting calculation explains how to recognize its $2$-cells geometrically.
\end{dcomment}

\begin{dcomment}{development}{p131}{Unbounded zigzags obstruct infinite products}
Grothendieck identifies the right difficulty in commuting localization with an infinite product: a morphism in a localization is represented by a finite word, and the lengths needed in different coordinates may be unbounded. The following construction turns that observation into a counterexample, even when the localizing classes are strongly saturated.

\textbf{Proposition (a product of finite fences).}\ours
For $n\geq1$, let $P_n$ be the finite poset with objects $0,1,\ldots,2n$ and nonidentity relations
\[
 0<1>2<3>\cdots<2n-1>2n.
\]
Let $W_n$ consist of every arrow. Then each $P_n[W_n^{-1}]$ is equivalent to the terminal category, but
\[
 \left(\prod_{n\geq1}P_n\right)
 \left[\left(\prod_{n\geq1}W_n\right)^{-1}\right]
 \longrightarrow \prod_{n\geq1}P_n[W_n^{-1}]
\]
is not full. More precisely, two objects $x=(x_n)$ and $y=(y_n)$ in the category on the left can be connected by a morphism if and only if $\sup_n|x_n-y_n|<\infty$.

\textit{Proof.} The only nonidentity arrows in $P_n$ join adjacent integers; there are no composable pairs of nonidentity arrows. Inverting these arrows gives the free groupoid on the undirected path with vertices $0,\ldots,2n$. Between any two vertices there is exactly one reduced path. Hence this groupoid is equivalent to the terminal groupoid.

Put $P=\prod_nP_n$. Since every arrow of $P$ is inverted, an arrow in its localization is represented by a finite zigzag in $P$. One step in such a zigzag changes the $n$th coordinate by at most one, for every $n$. A zigzag with $k$ steps therefore implies $|x_n-y_n|\leq k$ for all $n$. This proves necessity of the bound.

Conversely, suppose $|x_n-y_n|\leq K$ for all $n$. In each fence take the unique shortest undirected path from $x_n$ to $y_n$. Its arrow directions alternate. If its first arrow has the wrong orientation for a common alternating pattern beginning with a forward arrow, prepend an identity. After the path finishes, append identities. Every coordinate now has a zigzag of the same length $K+1$ and the same pattern of directions. Taking coordinatewise products gives a zigzag in $P$. This proves sufficiency.

Take $x_n=0$ and $y_n=2n$. There is no arrow between them in $P[W^{-1}]$, while their images have exactly one arrow between them in the product of the localized categories. Thus the comparison is not full. The localizing classes are all-arrow classes, so they are strongly saturated.\hfill$\square$

The original question for $\HOT_{W_\infty}$ nevertheless has a positive answer. Diagram categories of simplicial sets carry model structures with pointwise weak equivalences. For a disjoint union of shapes, the diagram model category is the product of the diagram model categories of the components. Its cofibrant-fibrant replacements are chosen coordinatewise. Between cofibrant-fibrant objects, every homotopy-category morphism has an actual representative, and homotopy can be tested coordinatewise using coordinatewise cylinder objects. Therefore the homotopy category of this product is the product of the homotopy categories, with no bound on the number of factors. The common model-category description provides the uniform presentations absent from the fence example; compare \citet[Proposition~1.30]{Groth2013}.

We checked the finite-word construction of localization in \citet[Chapter~I, \S1]{GabrielZisman1967} and the model-category derivator result just cited. We did not find this exact fence formulation in those sources. It is an elementary elaboration of the obstruction explicitly singled out in the manuscript, not a claim that the general phenomenon is new. In particular, it does not support replacing Grothendieck's tentative general argument with an assertion that all relative categories satisfy the coproduct axiom.
\end{dcomment}

\begin{dcomment}{exposition}{p133}{Constant coefficients explain the turn to homology}
Grothendieck turns from inverse image to homological direct image because the latter makes homotopy invariance of constant coefficients visible. For the space derivator, the homotopy colimit of a constant space $K$ over $S^{\mathrm{op}}$ is
\[
 \operatorname{hocolim}_{S^{\mathrm{op}}}K\simeq |N(S^{\mathrm{op}})|\times K
 \simeq |NS|\times K.
\]
One way to calculate it is the simplicial replacement: in degree $n$ it is a coproduct of copies of $K$, one for every string of $n$ composable arrows of $S^{\mathrm{op}}$. Realization therefore gives the displayed product. The reversal of a string identifies the realizations of the nerves of a category and its opposite.

For $f:S'\to S$, the canonical homological comparison runs from the expression for $S'$ to the expression for $S$ and is represented by
\[
 |Nf|\times\id_K:|NS'|\times K\longrightarrow|NS|\times K.
\]
When $f\in W_\infty$, this is a weak equivalence. The isomorphism written in the source may consequently be read in either direction after localization; before inversion the functorial map has the direction just displayed. For a general coefficient derivator the corresponding operation is tensoring with a homotopy type, rather than an internal Cartesian product of coefficients.

Adjunction now relates the calculation to the desired fully faithful restriction on constant objects. For constant coefficients $a,b$,
\[
 \Hom_{\mathbf D(S)}(p^*a,p^*b)
 \simeq\Hom_{\mathbf D(e)}(p_!p^*a,b).
\]
If the homological comparison is an isomorphism for every $a$, it induces a bijection on these Hom sets for every $b$. Conversely, bijectivity for every $b$ identifies the comparison by Yoneda. This is why the direct-image formalism gives exactly the statement Grothendieck wants, rather than merely an invariant of isomorphism classes.

The axiom numbers in these pages belong to the evolving manuscript. The condition called Der~3 here should not be confused with the usual modern axiom bearing that number, which asks for both Kan adjoints. Likewise the footnote retracting the unrestricted claim about Der~2 records a genuine limitation of the argument for arbitrary relative categories. The development at manuscript p.~130 gives a counterexample even when $W$ is strongly saturated, using a countably discrete diagram shape. It does not address the restriction of that question to finite shapes. The model-category theorem establishes the pointwise detection axiom for the space derivator; it does not vindicate the cancelled general assertion.
\end{dcomment}

\begin{dcomment}{exposition}{p136}{Changing the structure map gives direct image}
Grothendieck's proposed direct image has a particularly economical description: keep the total category and compose its map to the base. The substantial point is that this operation respects right-fine equivalence even though inverse image did not respect it on arbitrary representatives. The proof uses a class of test changes of base which is itself stable under pullback.

Let $u:X\to Y$ be in $W_{S'}^{\mathrm d}$ and let $f:S'\to S$. For an object $s\in S$, form
\[
 E=(s\backslash S)\times_S S'\longrightarrow S'.
\]
This projection is a discrete opfibration: an object is an arrow $s\to f(s')$, and an arrow out of $s'$ has a unique lift obtained by composing that arrow in $S$. Right-fine equivalences remain global $W$-equivalences after pullback along such maps. Applying this characterization to $u$ shows that
\[
 s\backslash(X/S)\longrightarrow s\backslash(Y/S)
\]
is in $W$. Since this holds for every $s$, the same underlying functor, regarded over $S$, belongs to $W_S^{\mathrm d}$.

The characterization can also be checked directly from the comma resolution. For a discrete opfibration $q:E\to S'$ and $e\in E$, unique lifting gives
\[
 e\backslash(X\times_{S'}E)\cong q(e)\backslash X.
\]
Thus the map between the comma resolutions over $E$ is a weak equivalence on fibres. The localizer's fibrewise-to-total property makes it a global weak equivalence, and the comparison from each comma resolution to its total category is a weak equivalence. Two-out-of-three gives the desired global conclusion. This uses only the comparison machinery already available.

The square of coslices considered in the source is Cartesian on morphisms as well as objects: an arrow is a pair of commuting triangles whose images in the bottom-right coslice agree. The universal property of the original pullback gives exactly one arrow in the top-left coslice with those data. Therefore the author's object check extends immediately to a proof of the asserted categorical pullback. The term ``isomorphisme local'' on ms.~136 must be read in its co-local sense, since it is outgoing coslices that occur throughout this argument.

We obtain
\[
 f_!q_{S'}(X)=q_S(X\xrightarrow{}S'\xrightarrow{f}S).
\]
Composition of structure maps gives the composition law for these functors. Their compatibility with $\rho$ and $\sigma$ is equally formal: the routes around the displayed square apply the same total-category functor and then invert a larger class. Neither observation says that direct image preserves locally constant objects. The interval example above shows that it need not.
\end{dcomment}

\begin{dcomment}{examples}{p137}{Integration and the local-constant reflection}
Grothendieck distinguishes two compatibility statements here because integration changes the support of a relative type. The inclusion $a:\{0\}\to S=(0\to1)$ already exhibits the phenomenon. The one-point object over $\{0\}$ is locally constant. Its direct image in the full relative category has comma fibres
\[
 (0\backslash\{0\})\simeq e,
 \qquad (1\backslash\{0\})=\varnothing.
\]
It is therefore not locally constant on $S$. Its reflection is the constant point, as calculated in the earlier block on ms.~128. Thus $a_!i_{\{0\}}$ and $i_Sa_!^{\mathrm{lc}}$ cannot be identified in general.

The commutation with reflection goes in the direction stated in the theorem:
\[
 \rho_S f_!\simeq f_!^{\mathrm{lc}}\rho_{S'}.
\]
There is a second explanation once the adjunctions have been established. The inclusions $i_S$ commute with restriction, and $\rho_S$ is left adjoint to $i_S$. Both functors in this display are consequently left adjoint to the same composite restriction followed by inclusion. The uniqueness of adjoints identifies them. In contrast, preserving the essential image of $i_{S'}$ under $f_!$ is a further property of $f$.

For ordinary weak equivalences, the Hot-fibration hypothesis in footnote~67 supplies precisely that property. If $s\to t$ is an arrow of $S$, the induced functor $t\backslash S'\to s\backslash S'$ is a weak equivalence. For a locally constant coefficient $F$, the pointwise value $(f_!F)(s)$ is the homotopy colimit of the pulled-back coefficient over $(s\backslash S')^{\mathrm{op}}$. A weak equivalence of indexing categories preserves this homotopy colimit with locally constant coefficients. Thus every transition of $f_!F$ is an equivalence: $f_!$ carries locally constant objects to locally constant objects, and $f_!i_{S'}\simeq i_S f_!^{\mathrm{lc}}$. This is the space-valued case of \citet[Proposition~1.24, applied to the opposite indexing convention]{Cisinski2009LocallyConstant}; the proof there makes the same pointwise homotopy-colimit comparison. It gives the affirmative answer sought in the footnote, while retaining the hypothesis that the preceding interval example shows to be substantive.

For the terminal map $p:S\to e$, composition of structure maps simply forgets the base. If a category $X\to S$ represents a coherent family of homotopy types, then $p_!q_S(X)$ is the homotopy type of the total category $X$. For a category of elements this is the homotopy-colimit formula expressed geometrically. It explains the word ``integration'': the coherent fibres, with their transport, are assembled into their total homotopy type. One should retain the word ``coherent'' here, since an ordinary diagram in the homotopy category alone does not determine that total type.
\end{dcomment}

\begin{dcomment}{exposition}{p139}{The unit comes from the comma replacement}
The unit must compare an object over $S'$ with the result of forgetting down to $S$ and then deriving inverse image. Forgetting the intermediate base can destroy the useful smoothness or fibration hypothesis. This is why $\Phi_S$ appears between those two operations, even if the original category was already fibred over $S'$.

It is helpful to describe the construction on objects. For $p:X\to S$, the category $\Phi_S(X)$ has objects $(s,x,v:s\to p(x))$ and is projected to $s$. A morphism from $(s,x,v)$ to $(t,y,w)$ is a pair $(h:s\to t,k:x\to y)$ satisfying $p(k)v=wh$. The unit $i_X$ sends $x$ to $(p(x),x,\id_{p(x)})$. Thus for $p':X'\to S'$ the unit of the desired adjunction is represented by
\[
 \alpha_{X'}:X'\longrightarrow\Phi_S(X'/S)\times_S S',
 \qquad x\longmapsto
 \bigl((f p'(x),x,\id_{f p'(x)}),p'(x)\bigr).
\]
The same formula on arrows uses $f p'(k)$ in the first coordinate and $p'(k)$ in the second. It is a functor over $S'$ and is natural in every functor $X'\to Y'$ over $S'$.

Equivalently, it is uniquely characterized by its two projections:
\[
 \pr_1\alpha_{X'}=i_{X'/S},
 \qquad \pr_2\alpha_{X'}=p'.
\]
This characterization is the ordinary slice adjunction between composition with $f$ and pullback along $f$, applied before localization. The target of $\alpha_{X'}$ represents $f^*f_!q_{S'}(X')$ because $\Phi_S(X'/S)$ is a split fibration over $S$. Hence the natural map descends to $\id\to f^*f_!$.

Naturality before localization is enough here, but it is important to use the correct base in the replacement. The top-right object in the source's square on ms.~139 is printed as $\Phi_S(X'/S')$. Its structure map and the preceding formula show that $\Phi_S(X'/S)$ is intended. This is a source slip, retained in the transcription and registered separately. No replacement over $S'$ can repair the loss of smoothness caused by composing its structure map with an arbitrary $f$.
\end{dcomment}

\begin{dcomment}{exposition}{p143}{The counit is a fraction over the original base}
Grothendieck's first formula for the counit is correct when the representative $X\to S$ is smooth: strict pullback represents derived inverse image, and projection
\[
 X\times_S S'\longrightarrow X
\]
represents $f_!f^*q_S(X)\to q_S(X)$. To check the first triangle identity, however, one needs the counit on $f_!q_{S'}(X')$, whose natural representative need not be smooth over $S$. The extension to arbitrary representatives is the small but essential issue being worked out on ms.~141--143.

For arbitrary $X$, put $R(X)=\Phi_S(X)\times_S S'$. The replacement map $i_X:X\to\Phi_S(X)$ is invertible in $\HOT_W(S)$, so the valid formula is
\[
 \beta_X=[i_X]^{-1}[\pr_1]:q_S(R(X))\longrightarrow q_S(X).
\]
Both arrows in this fraction are considered over $S$. Their naturality follows from the naturality of $i$ and of projection. If $X$ is smooth, the map $i_X\times_S S'$ is a right-fine equivalence over $S'$, and the natural square with the two projections identifies this general formula with the simple projection formula above.

The abandoned attempt instead inverts $i_X\times_S S'$ without a hypothesis making that pullback a weak equivalence. A different interval calculation shows the problem directly. Take $S=(0\to1)$, $X=\{1\}\to S$, and $S'=\{0\}$. The strict fibre of $X$ over $0$ is empty, while the fibre of $\Phi_S(X)$ is $0\backslash X\simeq e$. The pullback of $i_X$ is therefore $\varnothing\to e$. It is not a weak equivalence, either over $S'$ or after regarding it over $S$; its comma fibre at $0$ already fails. There is no inverse with which to form the proposed fraction.

The distinction is between an inverse in the localization over $S$, which is available, and the pullback of that inverse on arbitrary representatives, which is not represented by pulling back its denominator. The corrected roof leaves the denominator over $S$ and uses projection from a replacement on which inverse image can be computed. This is the same discipline that governs total derived functors: apply the operation to a suitable replacement, then use the comparison isomorphism in the category where it is known to exist.
\end{dcomment}

\begin{dcomment}{exposition}{p145}{The last square is naturality}
The argument has reached the last square before the second triangle identity. No new homotopical input is required. The remaining equality is the naturality of the comma-replacement map, and making that equality explicit finishes the adjunction on these pages.

For the first triangle, start with $X'\to S'$, regarded over $S$ after composition. The defining relation for the unit is
\[
 \pr_1\alpha_{X'}=i_{X'}.
\]
The general formula for the counit gives $\beta_{X'}=i_{X'}^{-1}\pr_1$ in the localization over $S$. Consequently $\beta_{X'}\alpha_{X'}=\id_{X'}$. This is the identity $\beta f_!\circ f_!\alpha=\id_{f_!}$, and its proof applies without any smoothness hypothesis on $X'$ over $S$.

For the second triangle we may represent $\xi$ by a fibration $X\to S$. Put $X'=X\times_S S'$ and $r=\pr_1:X'\to X$. Since $X$ is smooth, the counit at $\xi$ is represented by $r$. To calculate its derived inverse image on the possibly nonsmooth source $X'\to S$, replace that source by $\Phi_S(X')$. Functoriality of $\Phi_S$ then gives the relation
\[
 (i_X\times_S S')\,f^*(\beta_\xi)
 =\Phi_S(r)\times_S S'
\]
in the localization over $S'$. Here $i_X\times_S S'$ is invertible because both $X$ and $\Phi_S(X)$ are smooth over $S$.

It remains to prove
\[
 (\Phi_S(r)\times_S S')\alpha_{X'}=i_X\times_S S'.
\]
These are maps $X'\to\Phi_S(X)\times_S S'$. Their second projections both equal the structure map $X'\to S'$. Their first projections are respectively
\[
 \Phi_S(r)i_{X'}\qquad\text{and}\qquad i_Xr,
\]
which are equal by the naturality of $i:\id\to\Phi_S$. The universal property of the pullback proves the displayed equality already in $\Cat/S'$. Composing the preceding relation with $\alpha_{X'}$ and cancelling $i_X\times_S S'$ now gives
\[
 f^*(\beta_\xi)\alpha_{f^*\xi}=\id_{f^*\xi}.
\]
The two triangle identities establish $f_!\dashv f^*$.

There is a minor but consequential printed index in the intervening display on ms.~144. The map labelled $i_X\times_S\id_{S'}$ has target $\Phi_S(X)\times_S S'$, as in the next triangle, rather than $\Phi_S(X')\times_S S'$. The latter is the target of $\alpha_{X'}$ and the source of $f^*(\beta_\xi)$; identifying them would obscure the naturality square just used. The edition prints the prime, so it is preserved as a source slip. The proof given here completes this last square and makes no claim of originality.
\end{dcomment}


\begin{dcomment}{development}{p130}{Strong saturation does not imply pointwise detection}
Footnote~66 withdraws the assertion that every relative category defines a prederivator satisfying Der~2, even when its weak equivalences are strongly saturated. There is a distinction between two statements. If an actual natural transformation has invertible evaluations after localization, strong saturation makes each component a weak equivalence; that transformation is therefore inverted in the diagram category. An arbitrary morphism of the localized diagram category, however, need not have an actual natural transformation as representative. The following construction shows that this distinction can defeat pointwise detection. It also connects this footnote to the unbounded-length obstruction discussed on ms.~131.

\textbf{Proposition (a strongly saturated counterexample to Der~2).}\ours
There exist a small category $M$, a strongly saturated class $W$, and a countable discrete category $I$ such that
\[
 M^I[(W^I)^{-1}]\longrightarrow\prod_{i\in I}M[W^{-1}]
\]
is not conservative. In fact, the source has a morphism whose evaluations are all isomorphisms but has no morphism in the reverse direction. The category $M$ can be chosen to be a disjoint union of finite categories.

\textit{Proof: the finite categories.}
Let $R$ be the walking retraction, with objects $0,1$ and arrows
\[
 i:0\longrightarrow1,\qquad r:1\longrightarrow0,\qquad
 ri=\id_0,\qquad e=ir\ne\id_1.
\]
Its five arrows are $\id_0,\id_1,i,r,e$; in particular, its only isomorphisms are its identities. For example, this is the subcategory of sets generated by the inclusion of a singleton into a two-element set and its retraction.

Fix a positive odd integer $n$. Define a finite poset $Q_n$ with elements
\[
 x_0,\ldots,x_{n-1},\quad y_0,\ldots,y_{n-1},\quad z.
\]
For writing its order relations only, put $x_n=y_n=z$. Generate the order by the two fences
\[
 x_0\longleftarrow x_1\longrightarrow x_2\longleftarrow\cdots
 \longleftarrow x_n,
 \qquad
 y_0\longleftarrow y_1\longrightarrow y_2\longleftarrow\cdots
 \longleftarrow y_n,
\]
together with $x_j\leq y_j$ and $z\leq y_j$ for $0\leq j<n$. An arrow in these displays means an order relation from its source to its target. The generated relation is antisymmetric: no $y_j$ with $j<n$ reaches an $x_k$ or $z$, and the nonidentity arrows among the $x$'s have odd source and even target. Among the $x$'s the only nonidentity comparisons are the fence edges. In particular, $z=x_n$ reaches just $x_{n-1}$ among the other $x$'s. The full subposet on the $y$'s has initial object $y_n=z$.

Let $C_n$ be the full subcategory of the finite product category $R\times Q_n$ on the objects
\[
 X_j=(0,x_j),\quad Y_j=(1,y_j)\quad(0\leq j<n),\qquad
 X_n=(0,z),\quad Y_n=(1,z).
\]
Thus $X_n$ and $Y_n$ are distinct objects, despite having the same second coordinate. Write $F_n:C_n\to R$ for projection and set
\[
 W_n=F_n^{-1}(\operatorname{Iso}R).
\]
This class is strongly saturated. Indeed, $F_n$ inverts $W_n$ and hence factors through $C_n[W_n^{-1}]$. If an original arrow becomes invertible in that localization, its image under $F_n$ is an isomorphism, which is exactly the condition defining $W_n$. This argument also confirms that the section, the retraction and the idempotent $e$ do not inadvertently become weak equivalences.

\textit{A short morphism with a long inverse.}
There are weak arrows $c_j:Y_n\to Y_j$ given by $(\id_1,z\leq y_j)$, with $c_n=\id$, and arrows $i_j:X_j\to Y_j$ given by $(i,x_j\leq y_j)$. The product category also supplies $r_n:Y_n\to X_n$, with
\[
 r_ni_n=\id_{X_n}.
\]
For every fence edge $a:X_s\to X_t$, let $b:Y_s\to Y_t$ be the corresponding weak arrow. The following equations hold strictly:
\[
 i_ta=bi_s,\qquad bc_s=c_t.
\]
In the localization put $t_j=c_j^{-1}i_j:X_j\to Y_n$. These equations give $t_ta=t_s$. Let $p_n:X_0\to X_n$ be the morphism represented by the entire weak fence, traversed from index $0$ to index $n$. Since every fence edge is weak, $p_n$ is invertible. Transporting the preceding equation along its edges gives
\[
 i_np_n=t_0=c_0^{-1}i_0.
\]
Consequently the three-arrow word
\[
 f_n=r_nc_0^{-1}i_0:X_0\longrightarrow X_n
\]
equals $p_n$ and is an isomorphism.

Nevertheless every zigzag representing a morphism $X_n\to X_0$ has at least $n$ arrows. To prove this, assign the heights
\[
 h(X_j)=j,\qquad h(Y_j)=n\qquad(0\leq j\leq n).
\]
Every original forward arrow lowers height by at most one. For arrows between $X$ objects this follows from the fence description. Arrows from an $X$ to a $Y$ cannot lower it. Arrows between $Y$ objects preserve it. Finally, a $Y$ object can map to an $X$ object only from $Y_n$, and then the target is $X_n$ or $X_{n-1}$; these are exactly the comparisons in $Q_n$ with a $y$ source and an $x$ target. Such an arrow also lowers height by at most one.

A weak arrow has the same $R$ coordinate at its two ends. Between $Y$ objects it preserves height, and between $X$ objects its two indices differ by at most one. Its formal inverse therefore lowers height by at most one as well. Every letter in a localization zigzag, whether a forward arrow or an inverted weak arrow, lowers height by at most one. A word from height $n$ to height $0$ consequently has length at least $n$. This argument concerns every possible representative and allows all the extra endomorphisms supplied by fullness in $R\times Q_n$.

\textit{The countable diagram.}
Take
\[
 M=\coprod_{n\text{ positive odd}}C_n,\qquad
 W=\coprod_{n\text{ positive odd}}W_n,
\]
and let $I$ be the discrete category of positive odd integers. The projections $F_n$ together give $F:M\to R$, and $W=F^{-1}(\operatorname{Iso}R)$ is strongly saturated by the same argument. Define four $I$-diagrams by choosing, in the $n$th component, $X_0,Y_0,Y_n,X_n$, respectively. The families $i_0,c_0,r_n$ give a single diagram-category zigzag
\[
 \mathcal X\longrightarrow\mathcal U
 \xleftarrow{\,W^I\,}\mathcal V\longrightarrow\mathcal Z.
\]
Its localized morphism $f:\mathcal X\to\mathcal Z$ evaluates to $f_n$ and is therefore pointwise invertible.

Suppose that any morphism $\mathcal Z\to\mathcal X$ existed in $M^I[(W^I)^{-1}]$. Choose a finite zigzag representing it, of length $L$. At the coordinate indexed by any odd $n>L$, evaluation gives a zigzag $X_n\to X_0$ with at most $L$ letters. All its objects remain in $C_n$, since the disjoint union has no arrows between its components. The height bound gives $n\leq L$, a contradiction. There is thus no reverse morphism at all, and $f$ is not invertible.\hfill$\square$

The use of infinitely many coordinates is essential to this proof; it makes no assertion about the same question restricted to finite nondiscrete shapes. It refutes the unrestricted statement in footnote~66, including its strongly saturated version. It also strengthens the failure of the infinite-product comparison: that comparison can fail to reflect isomorphisms, not merely fail to be full.

For comparison, a precise sufficient hypothesis is available without a model structure. Suppose $W$ is strongly saturated and every morphism in $M^I[(W^I)^{-1}]$ has a representative
\[
 X\xleftarrow{w}X'\xrightarrow{a}Y'\xleftarrow{v}Y,
 \qquad w,v\in W^I.
\]
If its evaluations are isomorphisms, then every $a_i$ becomes invertible in $M[W^{-1}]$. Saturation gives $a_i\in W$, so the displayed word is invertible. A left or right calculus of fractions is a special case, allowing one of the outside arrows to be an identity. The same argument proves the Der~2 part of \citet[Proposition~4.2.8, pp.~85--86]{LagkasNikolos2018}; Remark~4.2.19, p.~95, explicitly discusses the role of two- and three-arrow descriptions in localization of prederivators.

The counterexample above was constructed for this commentary. The literature checks located the stated fraction criterion and the usual model-category results, but did not locate this particular finite-category construction. Its proof is included in full; the marking records our contribution to the present audit and does not assert historical priority for the general failure of localization to commute with infinite products.
\end{dcomment}

\begin{dcomment}{exposition}{p139}{The replacement must be taken over the new base}
Grothendieck now has a particularly economical candidate for direct image: a category over $S'$ is regarded as a category over $S$. The difficulty is to show that this operation has the expected universal property after localization. Its prospective right adjoint is not arbitrary strict pullback. It is pullback after a replacement which remembers arrows in the base. The next seven manuscript pages make that distinction precise enough to construct an adjunction, including its two triangle identities.

Write $Q_S=\Phi_S$ and $i_X:X\longrightarrow Q_SX$. For a category $p:X\longrightarrow S$, an object of $Q_SX$ is a triple $(s,x,a:s\longrightarrow p(x))$. A morphism from $(s,x,a)$ to $(t,y,b)$ consists of $v:s\longrightarrow t$ and $u:x\longrightarrow y$ satisfying $p(u)a=bv$. The projection to $s$ is the split fibration associated with the contravariant functor $s\longmapsto s\backslash X$. The map $i_X$ sends $x$ to $(p(x),x,\id_{p(x)})$. This recalls the comma replacement discussed at manuscript p.~20 (``Which comma construction repairs which variance''), but here its role is more demanding: it supplies the unit of a derived adjunction.

For $f:S'\longrightarrow S$, denote by $L$ the functor induced by composition with $f$, and by $R$ derived inverse image. If $[X]$ denotes $\hot_{S,W}(X)$, their formulas are
\[
 L[X']=[X'/S],\qquad R[X]=[Q_SX\times_SS'].
\]
The replacement in the second formula is over $S$. Even when $X'\longrightarrow S'$ is already a fibration, its composite with an arbitrary $f$ need not be one. Replacing over $S'$ would therefore solve the wrong problem. This is the reason for the otherwise rather insistent occurrences of $X'/S$ in the argument. The occurrence of $Q_S(X'/S')$ in the square on manuscript page 139 is a slip already present in the edition; its intended meaning is $Q_S(X'/S)$.

The unit can be written before localization. Its component is the unique $S'$-functor
\[
 a_{X'}:X'\longrightarrow Q_S(X'/S)\times_SS'
\]
whose first projection is $i_{X'/S}$ and whose second projection is the structure map of $X'$. Explicitly, it sends $x'$ to the triple with base object $f(p'(x'))$ and identity arrow, together with $p'(x')\in S'$. The construction is natural in $X'$. Passing through the localization therefore gives $\alpha:\id\longrightarrow RL$.

The counit is slightly different. If $X\longrightarrow S$ is smooth, the projection $X\times_SS'\longrightarrow X$ represents $\beta_{[X]}:LR[X]\longrightarrow[X]$. For an arbitrary representative, the formula is instead
\[
 \beta_{[X]}=[i_X]^{-1}
 [\pr_1:Q_SX\times_SS'\longrightarrow Q_SX].
\]
Both arrows in this fraction lie over $S$, and the arrow inverted is one already known to belong to $W_S^{\mathrm d}$. The forgetful functor $Q_SX\longrightarrow X$ on underlying categories does not generally lie over $S$: the two structure maps at $(s,x,a)$ have values $s$ and $p(x)$. The fraction is thus doing necessary work.

The formulas also explain why naturality survives the change of representatives. For an arrow $u:X\longrightarrow Y$ over $S$, the equalities $Q_S(u)i_X=i_Yu$ and $Q_S(u)\pr_1=\pr_1(Q_S(u)\times_SS')$ give the naturality square for $\beta$ after the $i$'s have been inverted. The construction uses the localization universal property, rather than a choice of inverses in the original category. This is precisely the amount of structure needed for the adjunction proof that follows.
\end{dcomment}

\begin{dcomment}{examples}{p142}{A pullback can destroy the replacement equivalence}
Grothendieck's first attempt to construct the counit asks whether the map $i_X$ can be inverted after an arbitrary pullback. The walking arrow shows exactly where that attempt breaks. We work with ordinary weak homotopy equivalences, $W=W_\infty$, so that an empty category cannot be equivalent to a point.

Let $S=[1]=(0\longrightarrow1)$ and let $X=\{1\}\longrightarrow S$. An object of $\Phi_SX$ is an arrow of $S$ with target $1$. There are two such objects, $0\longrightarrow1$ and $\id_1$, and one nonidentity morphism between them. The projection identifies $\Phi_SX$ with $S$ over itself. Consequently
\[
 i_X:\{1\}\longrightarrow S
\]
is right-fine: for $s=0$ and $s=1$, the functor $s\backslash X\longrightarrow s\backslash S$ is a weak equivalence between contractible categories.

Now pull back along $j:\{0\}\longrightarrow S$. The resulting map is
\[
 X\times_S\{0\}=\varnothing
 \longrightarrow
 \Phi_SX\times_S\{0\}=\{0\}.
\]
It is not in $W$, nor is it a right-fine equivalence over the point. Regarded again as a map over $S$, it still fails the right-fine test at $0$. Thus the proposed inverse of $\pi_X=i_X\times_S\id_{S'}$ need not exist in either localization relevant to the attempted argument.

The calculation also distinguishes the two inverse images without any abstract machinery. Strict pullback of $X$ to $0$ is empty, whereas derived pullback of $[X]$ to $0$ is a point. The missing object records the arrow $0\longrightarrow1$. A comma fibre retains that arrow; a strict fibre does not.

There is no defect in the replacement itself. The map $i_X$ becomes an isomorphism over $S$, exactly as intended. The invalid move was to apply a functor which had not been shown to preserve that class of equivalences on all of $\Cat/S$. On manuscript page 143 Grothendieck repairs the argument himself: he first projects to $\Phi_SX$, and then inverts $i_X$ in the localization over $S$. No pullback of an arbitrary denominator is required.
\end{dcomment}
